% concatenated from FP_Neural_OT_wo_Implicit_Differentiation_arxiv.tex, ex_shared_arxiv.tex
% SIAM Article Template
% \documentclass[review,onefignum,onetabnum]{siamonline250211}
\documentclass[onefignum,onetabnum]{siamonline250211}

% Information that is shared between the article and the supplement
% (title and author information, macros, packages, etc.) goes into
% ex_shared.tex. If there is no supplement, this file can be included
% directly.

% SIAM Shared Information Template
% This is information that is shared between the main document and any
% supplement. If no supplement is required, then this information can
% be included directly in the main document.


% Packages and macros go here
\usepackage{lipsum}
\usepackage{amsfonts}
\usepackage{graphicx}
\usepackage{epstopdf}
% \usepackage{algorithmic}
\ifpdf
  \DeclareGraphicsExtensions{.eps,.pdf,.png,.jpg}
\else
  \DeclareGraphicsExtensions{.eps}
\fi


\usepackage{amsmath}

\usepackage{amssymb}% http://ctan.org/pkg/amssymb
\usepackage{pifont}% http://ctan.org/pkg/pifont

\newcommand{\mycmark}{\green{\cmark}}
\newcommand{\myxmark}{\deepred{\xmark}}
\usepackage{booktabs}
\usepackage{subcaption}
\usepackage{tabularx}
\usepackage{multirow}
\usepackage{mathtools}
\usepackage{tikz}
\usetikzlibrary{positioning} 
\usetikzlibrary{calc}
% Prevent itemized lists from running into the left margin inside theorems and proofs
\usepackage{enumitem}
\setlist[enumerate]{leftmargin=.5in}
\setlist[itemize]{leftmargin=.5in}

\newtheorem{assumption}[theorem]{Assumption}
% Add a serial/Oxford comma by default.
\newcommand{\creflastconjunction}{, and~}

\DeclareMathOperator*{\argmin}{arg\,min}
\DeclareMathOperator*{\argmax}{arg\,max}

% Used for creating new theorem and remark environments
\newsiamremark{remark}{Remark}
\newsiamremark{hypothesis}{Hypothesis}
\crefname{hypothesis}{Hypothesis}{Hypotheses}
\newsiamthm{claim}{Claim}
\newsiamremark{fact}{Fact}
\crefname{fact}{Fact}{Facts}


% Sets running headers as well as PDF title and authors
\headers{Fixed-Point Neural Optimal Transport}{Y. Park, E. Gelphman, S. Osher and S. Wu Fung}

% Title. If the supplement option is on, then "Supplementary Material"
% is automatically inserted before the title.
\title{Implicit Neural Optimal Transport via Fixed-Point Optimization\thanks{Yesom Park and Eric Gelphman contributed to this work equally. 
\funding{This work was supported by DARPA HR00112590074, DoE DE-SC0026262, NSF 2208272, ARO W911NF241015, and NSF DMS 2309810.}}}


% Authors: full names plus addresses.
\author{Yesom Park\thanks{Department of Mathematics, University of California, Los Angeles (\email{yeisom@math.ucla.edu}), \email{sjo@math.ucla.edu}.} 
    \and Eric Gelphman \thanks{Department of Applied Mathematics and Statistics, Colorado School of Mines (\email{eric\_gelphman@mines.edu}), \email{swufung@mines.edu}.} 
\and Stanley Osher\footnotemark[2] 
    \and Samy Wu Fung\footnotemark[3]}

% \author{
% Yesom Park$^{\S}$\thanks{Department of Mathematics, University of California, Los Angeles, Los Angeles, CA 90095 (\email{yeisom@math.ucla.edu}, \email{sjo@math.ucla.edu}).}
% \and Eric Gelphman$^{\S}$\thanks{Department of Applied Mathematics and Statistics, Colorado School of Mines, Golden, CO 80401 (\email{eric\_gelphman@mines.edu}, \email{swufung@mines.edu}).}
% \and Stanley Osher\footnotemark[2]
% \and Samy Wu Fung$^{\P}$\footnotemark[3]
% }
% \begingroup
% \renewcommand{\thefootnote}{}
% \footnotetext{$^{\S}$ These authors contributed equally to this work.}
% \footnotetext{$^{\P}$ Corresponding author.}
% \endgroup


\usepackage{amsopn}
\DeclareMathOperator{\diag}{diag}


%%% Local Variables: 
%%% mode:latex
%%% TeX-master: "ex_article"
%%% End: 

\usepackage{algorithm}
\usepackage{algpseudocode}
% Optional PDF information
\ifpdf
\hypersetup{
  pdftitle={Fixed-Point Neural Optimal Transport without Implicit Differentiation},
  pdfauthor={D. Doe, P. T. Frank, and J. E. Smith}
}
\fi

% The next statement enables references to information in the
% supplement. See the xr-hyperref package for details.
% \usepackage{algorithm}
\usepackage{algpseudocode}
% \externaldocument[][nocite]{ex_supplement}

% FundRef data to be entered by SIAM
%<funding-group specific-use="FundRef">
%<award-group>
%<funding-source>
%<named-content content-type="funder-name"> 
%</named-content> 
%<named-content content-type="funder-identifier"> 
%</named-content>
%</funding-source>
%<award-id> </award-id>
%</award-group>
%</funding-group>

\begin{document}

\maketitle

% REQUIRED
\begin{abstract}
We propose an implicit neural formulation of optimal transport that eliminates adversarial min--max optimization and multi-network architectures commonly used in existing approaches. Our key idea is to parameterize a single potential in the Kantorovich dual and reformulate the associated c-transform as a proximal fixed-point problem. This yields a stable single-network framework in which dual feasibility is enforced exactly through proximal optimality conditions rather than adversarial training.
Despite the inner fixed-point computation, gradients can be computed without differentiating through the fixed-point iterations, enabling efficient training without requiring implicit differentiation. We further establish convergence of stochastic gradient descent. The resulting framework is efficient, scalable, and broadly applicable: it simultaneously recovers forward and backward transport maps and naturally extends to class-conditional settings.
Experiments on high-dimensional Gaussian benchmarks, physical datasets, and image translation tasks demonstrate strong transport accuracy together with improved training stability and favorable computational and memory efficiency.
\end{abstract}

% REQUIRED
\begin{keywords}
Optimal Transport, Fixed-point Iteration, Deep Learning, Kantorovich Dual, Convergence 
\end{keywords}

% REQUIRED
\begin{MSCcodes}
49Q22, 68T07, 65K10
\end{MSCcodes}

\section{Introduction}

Optimal transport (OT) is a fundamental problem of finding a mapping between probability distributions that minimizes a prescribed transportation cost. Owing to its solid theoretical foundation and wide applicability, OT has been successfully employed in diverse fields such as traffic control \cite{carlier2008optimal,danila2006optimal, barthelemy2006optimal}, biomedical data analysis \cite{schiebinger2019optimal, koshizuka2022neural, bunne2023learning}, generative modeling \cite{wang2021deep,onken2021ot,zhang2023mean, liu2019wasserstein}, and domain adaptation \cite{courty2016optimal, courty2017joint, damodaran2018deepjdot, balaji2020robust}.

Recent advances in deep learning have led to a surge of interest in scalable OT solvers based on neural parameterizations. Early approaches are rooted in the Monge formulation \cite{lu2020large, xie2019scalable} and its relaxation to the Kantorovich framework \cite{makkuva2020optimal}. Despite their theoretical rigor, these formulations often entail significant computational complexity, particularly in high-dimensional settings. A predominant line of work is based on the Kantorovich dual formulation, which recasts the OT problem as a saddle-point optimization over transport maps and dual potentials \cite{liu2019wasserstein, taghvaei20192, korotin2021neural, liu2021learning, choi2024improving}. While this perspective enables scalable algorithmic implementations, it typically necessitates adversarial training of multiple neural networks, thereby introducing optimization instability, sensitivity to hyperparameter selection, and convergence difficulties. These challenges are further exacerbated in WGAN-based methods \cite{arjovsky2017wasserstein, liu2019wasserstein}, particularly as the dimensionality of the problem increases.

To leverage additional structural properties, several studies focus on the quadratic-cost setting, where Brenier’s theorem \cite{benamou2000computational} guarantees that the OT map can be expressed as the gradient of a convex potential. This observation has motivated the use of input convex neural networks (ICNNs) \cite{amos2017input} and related architectures \cite{taghvaei20192, korotin2021neural}. In this direction, prior work typically parameterizes either the primal and dual potentials separately~\cite{makkuva2020optimal, taghvaei20192}, or directly models the convex conjugate minimizer ~\cite{amos2022amortizing}. In addition, weak formulations have been explored to directly parameterize transport maps \cite{backhoff2019existence, korotin2022neural, asadulaev2022neural}. Nevertheless, many of these approaches continue to rely on auxiliary networks, alternating optimization procedures, or adversarial training, and thus inherit the limitations associated with min–max optimization.

An alternative line of research formulates OT as a dynamical system via continuous flows \cite{yang2020potential, tong2020trajectorynet, onken2021ot, huguet2022manifold}. These methods typically require solving ordinary differential equations (ODEs) or stochastic differential equations (SDEs), which imposes considerable computational overhead during both training and inference. Regularized variants, including entropic and $f$-divergence-based formulations \cite{genevay2016stochastic, seguy2017large, daniels2021score, gushchin2023entropic}, can improve numerical stability but generally introduce bias, leading to transport maps that deviate from the true OT solution.

Various efforts have been made to mitigate these challenges through improved optimization strategies, such as natural gradient methods \cite{shen2020sinkhorn, liu2024natural}, as well as regularization techniques including $L^2$ penalties \cite{genevay2016stochastic, sanjabi2018convergence} and cycle-consistency constraints \cite{genevay2016stochastic, sanjabi2018convergence, korotin2019wasserstein, korotin2021continuous}. However, these approaches do not fundamentally resolve the reliance on complex optimization schemes or multiple model components \cite{korotin2021neural, fan2021variational}.
More recently, formulations based on Hamilton–Jacobi–Bellman (HJB) equations \cite{park2025implicit} have been proposed to eliminate adversarial training by casting OT as a single-objective optimization problem \cite{park2025neural}. In particular, characteristic-based representations enable the recovery of both forward and backward transport maps within a unified framework. However, such methods require solving the HJB equation over the entire computational domain, which significantly limits their scalability in high-dimensional settings.



{\color{black}
In this work, we build upon single-potential formulation on quadratic OT \cite{taghvaei20192} and investigate an implicit neural framework that recovers both forward and backward transport maps through a single minimization problem defined over one neural network. 
Starting from the Kantorovich dual problem, we parameterize the dual potential as the negative of a convex function represented by a neural network $g_\theta$. Under this parameterization, the associated $c$-transform reduces to a strongly convex proximal optimization problem, whose minimizer can be computed efficiently via standard fixed-point iterations. This a single-objective optimization framework in which dual feasibility is enforced directly through the optimality conditions of the proximal operator, thereby eliminating the need for saddle-point formulations, auxiliary conjugate networks, and alternating optimization procedures. 

A central contribution of this work is the analysis of stochastic optimization under inexact inner fixed-point iteration. A fundamental challenge in implicit single-minimization OT formulations is that the convergence behavior of the model depends critically on the accuracy of the inner solver used to evaluate the $c$-transform. In practice, however, the associated inner problem cannot be solved exactly and must instead be approximated using a finite number of iterations.
We analyze the approximate stochastic gradient induced by this inexact fixed-point computation and establish convergence guarantees for the resulting SGD dynamics. In particular, we show that if the fixed-point residual is controlled within a prescribed tolerance, then the resulting approximate gradient remains a descent direction for the objective. Under standard diminishing step-size conditions, we further prove that convergence of SGD to a stationary point despite inexact inner solves. These results provide a theoretical justification for practical training with fixed-point iterations and remove the need for exact inner optimization in large-scale implicit neural OT.

We further extend the framework to class-conditional OT, where the objective is to learn transport maps that preserve class labels while minimizing transportation cost. We validate the proposed approach across a broad range of experimental settings. In particular, we demonstrate accurate recovery of transport maps on high-dimensional Gaussian benchmarks and strong performance on physics-based datasets involving complex structured distributions. We further evaluate the method on class-conditional transport tasks, including image translation experiments, where it achieves competitive performance while maintaining stable training dynamics and favorable computational and memory efficiency as the problem dimension increases.

The remainder of this paper is organized as follows. In Section \ref{sec:background}, we review the background on OT and related formulations. Section \ref{sec:method} introduces the implicit neural OT framework together with the convergence analysis of the resulting optimization procedure. In Section~\ref{sec:experiments}, we present experimental results across a range of benchmark datasets, including CCOT tasks and model ablation studies. Finally, Section \ref{sec:conclusion} concludes the paper and discusses future directions.}

% In this work, \black{we build upon single-potential formulation on quadratic OT} and investigate an implicit neural framework that recovers both forward and backward transport maps through a single minimization problem defined over one neural network. 
% % Specifically, we show that quadratic OT admits a significantly simpler and more direct formulation than those employed in existing neural approaches. 
% Starting from the Kantorovich dual problem, we parameterize the dual potential as the negative of a convex function represented by a neural network $g_\theta$. Under this parameterization, the associated $c$-transform reduces to a strongly convex proximal optimization problem, whose minimizer can be computed efficiently via standard fixed-point iterations. This leads to a single-objective optimization problem in which dual feasibility is enforced exactly through the optimality conditions of the proximal operator, thereby eliminating the need for saddle-point formulations, auxiliary conjugate networks, and alternating optimization procedures. As a result, the entire model is defined by a single convex potential learned through a standard minimization framework.
% A key observation is that the gradient of the resulting objective does not require differentiation through the inner minimization. By exploiting the first-order optimality condition, the dependence of the objective on the proximal minimizer simplifies, yielding a closed-form expression for the gradient. Consequently, training can be carried out using standard backpropagation applied solely to the potential network, without implicit differentiation or unrolled optimization. These properties lead to a method that is both computationally efficient and stable in practice: removing adversarial training mitigates instability associated with min–max optimization, while avoiding differential equation solvers eliminates the overhead inherent in dynamical formulations. Furthermore, fixed-point iterations enable fast and scalable computation across dimensions.

% The proposed framework relies on a single network architecture to recover both forward and backward OT maps, reducing model complexity and improving computational efficiency. In addition, the formulation naturally extends to CCOT problems. Empirically, we demonstrate that the proposed method accurately recovers transport maps on high-dimensional Gaussian benchmarks and remains effective on physics-based datasets, capturing complex and realistic distributions. The method also performs competitively on class-conditional transport tasks, including experiments on image data, while exhibiting stable training dynamics and favorable computational and memory efficiency as the dimensionality increases.

% \paragraph{Contributions}
% Our main contributions are summarized as follows:
% \begin{itemize}
%     \item We derive a single-network reformulation of quadratic OT from the Kantorovich dual, in which evaluating the objective reduces to solving a fixed-point problem.
%     \item We show that this objective can be differentiated without implicit differentiation, due to an exact cancellation induced by the first-order optimality conditions of the proximal operator.
%     \item We establish convergence of the resulting training procedure to a stationary point under inexact fixed-point evaluations, accounting for the fact that the fixed point is not solved exactly in practice.
%     \item We propose a single-network method that simultaneously recovers both forward and backward OT maps.
%     \item We validate our approach empirically, demonstrating strong performance and favorable scaling in high-dimensional settings, including class-conditional scenarios.
% \end{itemize}

% OT (OT) provides a principled framework for comparing probability distributions and constructing transport maps between them. With its strong theoretical foundation and broad practical relevance, OT has been widely applied in diverse areas, including traffic control, biomedical data analysis.

% In recent years, neural OT methods have become important tools for high-dimensional generative modeling, domain adaptation, and distribution alignment. However, most existing neural OT solvers rely on min--max formulations of the Kantorovich dual, which require the simultaneous optimization of multiple networks and often exhibiting instability, saddle-point dynamics, and degraded performance as dimension increases.

% For the quadratic cost, OT admits additional structure. Brenier's theorem guarantees that, under mild regularity conditions, the OT map can be expressed as the gradient of a convex potential. Existing neural approaches leverage this structure either by parameterizing the transport map directly or by alternating between dual potentials and their convex conjugates in adversarial training schemes. These formulations introduce auxiliary networks, saddle-point optimization, and approximate enforcement of dual constraints.

% In this work, we show that quadratic OT admits a significantly simpler formulation. Starting directly from the Kantorovich dual, we parameterize the dual potential as the negative of a convex neural network $g_\theta$. The associated $c$-transform reduces to a strongly convex proximal problem whose minimizer can be computed efficiently using standard first-order methods. This leads to a constrained optimization formulation in which dual feasibility is enforced exactly through the proximal optimality condition. As a consequence, training reduces to optimizing a single convex potential without min--max dynamics, auxiliary conjugate networks, or implicit differentiation.

% A key observation is that the gradient of the resulting objective does not require differentiating through the inner minimization. By first-order optimality of the proximal problem, all dependence on the derivative of the minimizer cancels exactly, yielding a simple closed-form gradient expression. Thus, quadratic OT can be trained using standard backpropagation applied only to the potential network itself.

% Empirically, we demonstrate that this formulation yields highly accurate transport maps on high-dimensional Gaussian benchmarks and performs competitively on class-conditional transport tasks, while maintaining stable computational and memory costs across dimensions.

% \paragraph{Contributions}
% Our main contributions are summarized as follows:
% \begin{itemize}
%     \item We derive a single-network reformulation of quadratic OT from the Kantorovich dual, in which evaluating the objective reduces to solving a fixed-point problem.
%     \item We show that this objective can be differentiated without implicit differentiation, due to an exact cancellation induced by the first-order optimality conditions of the proximal operator.
%     \item Recognizing that the fixed point is not solved exactly in practice, we establish convergence of the resulting training procedure to a stationary point under inexact fixed-point evaluations.
%     \item We validate our approach empirically, demonstrating strong performance and favorable scaling in high-dimensional settings.
% \end{itemize}

\section{Background: Quadratic OT and Convex Potentials}\label{sec:background}

\subsection{Wasserstein--2 Distance and Kantorovich Duality}

Let $\mu$ and $\nu$ be probability measures on $\mathbb{R}^d$ with finite second moments. For the quadratic cost
\begin{equation}
c(x,z) = \frac{1}{2}\|x - z\|^2,
\label{eq:quadratic_cost}
\end{equation}
the squared $2$-Wasserstein distance is defined as
\begin{equation}
W_2^2(\mu,\nu)
=
\inf_{\gamma \in \Pi(\mu,\nu)}
\int_{\mathbb{R}^d \times \mathbb{R}^d}
\frac{1}{2}\|x - z\|^2 \, d\gamma(x,z),
\end{equation}
where $\Pi(\mu,\nu)$ denotes the set of couplings with marginals $\mu$ and $\nu$.

The Kantorovich dual formulation states that
\begin{equation}
W_2^2(\mu,\nu)
=
\sup_{\varphi}
\left\{
\mathbb{E}_{x\sim\mu}[\varphi(x)]
+
\mathbb{E}_{z\sim\nu}[\varphi^c(z)]
\right\},
\end{equation}
where the $c$-transform is defined by
\begin{equation}
\varphi^c(z)
=
\inf_{y \in \mathbb{R}^d}
\left\{
\frac{1}{2}\|y - z\|^2 - \varphi(y)
\right\}.
\end{equation}

\subsection{Brenier's Theorem and Convex Potentials}

For the quadratic cost, OT admits additional structure. 
Brenier's theorem states that if $\mu$ is absolutely continuous, then there exists a convex function $u : \mathbb{R}^d \to \mathbb{R}$ such that the OT map from $\mu$ to $\nu$ is given by
\begin{equation}
T_{\mu\to\nu}(x) = \nabla u(x).
\end{equation}
Equivalently, writing $u(x)=\frac12\|x\|^2 + g(x)$ with $g$ convex, the map can be expressed as
\begin{equation}
T_{\mu\to\nu}(x) = x + \nabla g(x), 
\label{eq:T}
\end{equation}
which is a result of the following theorem. 
%\textcolor{red}{(EG) I think this is just $T_{\mu \rightarrow \nu}(x) = \nabla g(x)$ See proof in supplementary materials for why I think this is.}

\begin{theorem}
Suppose $\mu$ is absolutely continuous with respect to the Lebesgue measure on $\mathbb{R}^d$. Then, $\exists$ a unique OT map $T: \mathbb{R}^d \rightarrow \mathbb{R}^d$ with respect to the quadratic cost \eqref{eq:quadratic_cost} where $T_{\#}\mu = \nu$ and $T$ is given by \eqref{eq:T}.
\label{theorem:optimal_transport_map}
\end{theorem}
This is a standard result from~\cite[Theorem 10.28]{villani2008optimal} and we provide a proof in the appendix for completeness.

This characterization motivates parameterizing the dual potential as
\begin{equation}
\varphi(x) = -g(x),
\end{equation}
where $g$ is convex. Under this choice, the $c$-transform becomes
\begin{equation}\label{eq:c_transform}
(-g)^c(z)
=
\inf_{y \in \mathbb{R}^d}
\left\{
\frac{1}{2}\|y - z\|^2 + g(y)
\right\}.
\end{equation}
This expression is precisely the Moreau envelope (or quadratic inf-convolution) of $g$~\cite{tibshirani2025laplace, osher2023hamilton} and corresponds to the solution of a Hamilton-Jacobi equation~\cite{darbon2021connecting, darbon2021bayesian, heaton2024global, di2026operator}.

\subsection{Proximal Characterization and Transport Maps}

For convex and differentiable $g$, the minimizer
\begin{equation}
y^*(z)
=
\arg\min_y
\left\{
\frac{1}{2}\|y - z\|^2 + g(y)
\right\}
\end{equation}
is uniquely defined~\cite{parikh2014proximal}. 
% First-order optimality yields
% \begin{equation}
% 0
% =
% \nabla g(y^*(z)) - (z - y^*(z)),
% \qquad
% \text{or equivalently}
% \qquad
% y^*(z)
% =
% z - \nabla g(y^*(z)).
% \end{equation}
% Thus, $y^*(z)$ is characterized as the fixed point of a nonlinear operator.
The forward and backward transport maps admit the representations
\begin{equation}
T_{\mu\to\nu}(x)
=
x + \nabla g(x),
\qquad
T_{\nu\to\mu}(z)
=
y^*(z).
\end{equation}
Therefore, the quadratic OT problem can be formulated entirely in terms of a single convex potential $g$, whose $c$-transform is evaluated by solving a strongly convex minimization problem. This proximal characterization forms the basis of the implicit neural formulation developed in the next section.

\section{Implicit Neural OT}\label{sec:method}

We now present an implicit neural formulation of quadratic OT based on the proximal characterization of the $c$-transform.

\subsection{Parameterized Dual Formulation}

Recall the Kantorovich dual formulation for the quadratic cost:
\begin{equation}
W_2^2(\mu,\nu)
=
\sup_{\varphi}
\left\{
\mathbb{E}_{x\sim\mu}[\varphi(x)]
+
\mathbb{E}_{z\sim\nu}[\varphi^c(z)]
\right\}.
\end{equation}
We parameterize the dual potential as
\begin{equation}
\varphi_\theta(x) = -g_\theta(x),
\end{equation}
where $g_\theta : \mathbb{R}^d \to \mathbb{R}$ is assumed to be convex in its argument. Under this parameterization, the $c$-transform becomes
\begin{equation}
(-g_\theta)^c(z)
=
\inf_{y \in \mathbb{R}^d}
\left\{
\frac{1}{2}\|y - z\|^2 + g_\theta(y)
\right\}.
\end{equation}
Substituting into the dual objective yields
\begin{equation}
W_2^2(\mu,\nu)
=
\sup_\theta
\left\{
-
\mathbb{E}_{x\sim\mu}[g_\theta(x)]
+
\mathbb{E}_{z\sim\nu}
\left[
\inf_y
\left(
\frac{1}{2}\|y - z\|^2 + g_\theta(y)
\right)
\right]
\right\}.
\end{equation}
Equivalently, we minimize the negative dual objective:
\begin{equation}
\label{eq:dual_loss}
\mathcal{L}(\theta)
=
\mathbb{E}_{x\sim\mu}[g_\theta(x)]
-
\mathbb{E}_{z\sim\nu}
\left[
\inf_{y}
\left(
\frac{1}{2}\|y - z\|^2 + g_\theta(y)
\right)
\right].
\end{equation}

\subsection{Constrained Implicit Formulation}

For each $z \in \mathbb{R}^d$, define the proximal minimizer
\begin{equation}
\label{eq:prox_def}
y_\theta^\star(z)
=
\arg\min_{y \in \mathbb{R}^d}
\left\{
\frac{1}{2}\|y - z\|^2 + g_\theta(y)
\right\}.
\end{equation}

Since the objective is strongly convex in $y$, the minimizer exists and is unique. The $c$-transform can therefore be written explicitly as
\begin{equation}
(-g_\theta)^c(z)
=
\frac{1}{2}\|y_\theta^\star(z) - z\|^2
+
g_\theta\big(y_\theta^\star(z)\big).
\end{equation}
Substituting this expression into \eqref{eq:dual_loss}, the training problem becomes the constrained training problem given by
\begin{equation}
\label{eq:constrained_problem}
\begin{aligned}
\min_\theta \quad 
& 
\mathbb{E}_{x\sim\mu}[g_\theta(x)]
-
\mathbb{E}_{z\sim\nu}
\left[
\frac{1}{2}\|y_\theta^\star(z) - z\|^2
+
g_\theta\big(y_\theta^\star(z)\big)
\right] \\
\text{s.t.} \quad
&
y_\theta^\star(z)
=
\arg\min_{y}
\left\{
\frac{1}{2}\|y - z\|^2 + g_\theta(y)
\right\}
\quad
\text{for } z \sim \nu.
\end{aligned}
\end{equation}
This formulation is closely related to single-potential dual representations of quadratic OT~\cite{taghvaei20192}. Our contribution is to leverage this structure to develop a scalable and practically effective method that performs well across a wide range of datasets, together with a convergence analysis of the resulting optimization procedure under inexact evaluation of $y_\theta^\star(z)$.

In this formulation, learning reduces to optimizing a single convex potential $g_\theta$, with the model output implicitly defined by the proximal optimality condition that characterizes $y_\theta^\star(z)$. Importantly, dual feasibility is enforced by construction through~\eqref{eq:prox_def}, eliminating the need for adversarial min--max formulations with auxiliary networks~\cite{korotin2019wasserstein, korotin2022neural, heaton2022wasserstein, lin2021alternating}, as well as time-stepping approaches that require learning entire trajectories~\cite{onken2021ot, ruthotto2020machine}.

% This perspective places our approach within the framework of implicit deep learning~\cite{el2021implicit, fung2026generalization}, where the model output is defined implicitly via $y_\theta^\star$. However, as we show in Section~\ref{subsec:gradient_computation}, the structure induced by the proximal operator allows us to avoid implicit differentiation, even though computing $y_\theta^\star$ still requires solving a fixed-point problem.

% \subsection{Transport Maps}

Once $g_\theta$ is trained, the forward and backward transport maps are given by
\begin{equation}
T_{\mu\to\nu}(x)
=
x + \nabla g_\theta(x),
\qquad
T_{\nu\to\mu}(z)
=
y_\theta^\star(z).
\end{equation}

Both transport directions are therefore represented using a single convex potential. The backward map is computed as the unique minimizer of a strongly convex objective, while the forward map follows directly from Brenier's theorem. The overall training procedure, including the gradient computation described in the following subsection, is summarized in Algorithm~\ref{alg:implicit_ot}.

\begin{algorithm}[t]
\caption{Implicit Neural OT via Fixed-Point Proximal Updates} \label{alg:implicit_ot}
\begin{algorithmic}[1]
\Require Datasets $\mathcal{D}_\mu, \mathcal{D}_\nu$, initial parameters $\theta$, step size $\eta$, fixed-point steps $K$
\For{each training iteration}
    
    \State Sample minibatches $\{x_i\}_{i=1}^B \sim \mathcal{D}_\mu$, $\{z_i\}_{i=1}^B \sim \mathcal{D}_\nu$

    \State \textbf{// Fixed-point iteration for proximal operator}
    \For{each $z_i$ in minibatch}
        \State Initialize $y^{(0)} \gets z_i$
        \For{$k = 0, \dots, K-1$}
            \State $y^{(k+1)} \gets y^{(k)} - \alpha \big( \nabla_y g_\theta(y^{(k)}) + y^{(k)} - z_i \big)$
        \EndFor
        \State $\tilde{y}_i \gets \texttt{stop\_gradient}(y^{(K)})$ // Detach computational graph
    \EndFor

    \State \textbf{// Compute Loss}
    \State $g \gets \frac{1}{B} \sum_{i=1}^B \nabla_\theta g_\theta(x_i) - \frac{1}{B} \sum_{i=1}^B \nabla_\theta g_\theta(\tilde{y}_i)$

    \State \textbf{// Parameter update}
    \State $\theta \gets \theta - \eta \, g$

\EndFor
\end{algorithmic}
\end{algorithm}

\begin{remark}[PDE interpretation]
The proposed formulation admits an alternative interpretation from the perspective of partial differential equations (PDEs). In particular, the optimality condition of quadratic OT can be expressed in terms of the Hamilton--Jacobi (HJ) equation with a quadratic Hamiltonian. Owing to convexity, its solution can be characterized by the Hopf--Lax formula, which coincides with the $c$-transform in \eqref{eq:prox_def}. From this viewpoint, our method can be interpreted as solving the HJ equation via fixed-point iterations applied to the Hopf--Lax operator \cite{park2025fixed}. 

In contrast to prior works, which solve the HJ equation over the entire spatio-temporal computational domain and are therefore computationally demanding in high dimensions, we instead evaluate a local fixed-point map that enforces the corresponding optimality condition. This provides a causality-free mechanism for enforcing the PDE optimality condition without requiring the propagation of full PDE dynamics. Consequently, the proposed approach yields a significantly more efficient procedure while still enforcing the underlying optimality condition associated with the OT problem.
\end{remark}

\subsection{Gradient Computation Without Implicit Differentiation}
\label{subsec:gradient_computation}
{\color{black}
Since the model output is defined implicitly through a proximal optimization layer, one might expect that gradient computation requires implicit differentiation techniques as in implicit deep learning frameworks~\cite{el2021implicit, fung2022jfb, gelphman2025end, gelphman2026convergence}. However, the structure of formulation~\eqref{eq:constrained_problem} allows us to compute the gradient of the objective $\mathcal{L}$ without differentiating through the inner minimization that defines $y_\theta^\star(z)$. This cancellation-type property has also been observed in prior works~\cite{amos2022amortizing, taghvaei20192}, where the structure of the convex conjugate eliminates the need for implicit differentiation. Below, we show the formula for such derivative and provide its proof for completeness.}
% A key property of the formulation \eqref{eq:constrained_problem} is that differentiation does not require backpropagating through the inner minimization defining $y_\theta^\star(z)$. The gradient of $\mathcal{L}$ admits a remarkably simple expression.

\begin{lemma}
\label{lem:no_implicit_diff}
Assume $0 < \gamma < 1$, $g_\theta(y)$ is $\gamma$-weakly convex, continuously differentiable in $y$, and L-smooth in $\theta$, and that $y_\theta^\star(z)$ is the unique minimizer of
\begin{equation}
\min_y \left\{ \frac{1}{2}\|y - z\|^2 + g_\theta(y) \right\}.
\end{equation}
Then the gradient of $\mathcal{L}$ is given by
\begin{equation}
\label{eq:gradient_formula}
\frac{d\mathcal{L}}{d\theta}
=
\mathbb{E}_{x \sim \mu}
\left[
\frac{\partial g_\theta(x)}{\partial \theta}
\right]
-
\mathbb{E}_{z \sim \nu}
\left[
\frac{\partial g_\theta\big(y^\star_\theta(z)\big)}{\partial \theta}
\right].
\end{equation}
\end{lemma}

\begin{proof}
The forward pass of the network can be characterized as finding the fixed point of the operator
\begin{equation}
F_{\theta}(y) = y - \alpha(\nabla g + y - z) 
\end{equation}
% Applying the implicit function theorem like in the original JFB paper, we arrive at the expression for the true gradient
% \begin{equation}
% \frac{d y^*}{d \theta} = \left(I - \frac{\partial F_{\theta}}{\partial y}\right)^{-1}\frac{\partial F_{\theta}}{\partial \theta}
% \end{equation}
% The JFB approximation is
% \begin{equation}
% \frac{d y^*}{d \theta} \approx  \frac{\partial F_{\theta}}{\partial
% \theta} 
% \end{equation}
Assume that we can interchange $\mathbb{E}_{x}[\cdot]$ and $\mathbb{E}_{z}[\cdot]$ for any $x,z$. Then,
\begin{align}
\frac{d L}{d \theta} &= \mathbb{E}_{x \sim \mu}\left[\frac{d g}{d \theta}\right] - \mathbb{E}_{z \sim \nu}\left[\frac{d}{d \theta}(\frac{1}{2}\|y^* - z\|^2 + g(y^*))\right] \\
&=  \mathbb{E}_{x \sim \mu}\left[\frac{d g}{d \theta}\right] - \mathbb{E}_{z \sim \nu}\left[(y^* - z)\frac{d y^*}{d \theta} + \frac{ \partial g }{\partial y}\frac{d y^*}{d \theta} + \frac{\partial g}{\partial \theta}\right] \\
&= \mathbb{E}_{x \sim \mu}\left[\frac{\partial g}{\partial \theta}\right] - \mathbb{E}_{z \sim \nu}\left[\left((y^* - z) + \frac{ \partial g }{\partial y}\right)\frac{d y^*}{d \theta} + \frac{\partial g}{\partial \theta}\right].
\end{align}
% The JFB approximation to $\frac{d L}{d \theta}$, denoted $d_{\theta}^{JFB}$, is given by
% \begin{equation}
% d_{\theta}^{JFB} = \mathbb{E}_{x \sim \mu}\left[\frac{\partial g}{\partial \theta}\right] - \mathbb{E}_{z \sim \nu}\left[\left((y^* - z) + \frac{ \partial g }{\partial y}\right)\frac{\partial F_{\theta}}{\partial \theta} + \frac{\partial g}{\partial \theta}\right] 
% \end{equation}
If $y^*$ is a minimizer of $\frac{1}{2}\|y-z\|^2 + g_{\theta}(y)$, then
\begin{equation*}
\left( \nabla_{y}\left(\frac{1}{2}\|y-z\|^2 + g_{\theta}(y)\right)\right|_{y=y^*} = \left(2(\frac{1}{2}(y-z)) + \frac{\partial g(y)}{\partial y}\right|_{y=y^*} = y^* - z + \frac{\partial g(y^*)}{\partial y} = 0.
\end{equation*}
Thus, if the fixed point problem is solved exactly,
\begin{equation}
\nabla_{\theta}L = \mathbb{E}_{x \sim \mu}\left[\frac{\partial g(x)}{\partial \theta}\right] - \mathbb{E}_{z \sim \nu}\left[\frac{\partial g(y^*)}{\partial \theta}\right].
\end{equation}
\end{proof}

Lemma~\ref{lem:no_implicit_diff} shows that the gradient of the dual objective depends only on explicit derivatives of $g_\theta$, evaluated at samples from $\mu$ and at the proximal points $y_\theta^\star(z)$. In particular, no implicit differentiation, i.e., computation of $\frac{dy_\theta^\star}{d\theta}$ through the inner optimization problem, is required. This cancellation follows directly from the first-order optimality condition of the proximal operator and can be interpreted as an application of the envelope theorem. Consequently, backpropagation can be performed without differentiating through the inner fixed-point iterations used to compute $y_\theta^\star(z)$, making gradient computation computationally efficient in practice.

% While one might think that using a network to parameterize $g_\theta$ would require training an implicit network~\cite{el2021implicit, fung2022jfb, gelphman2025end, gelphman2026convergence}, 
% Lemma~\ref{lem:no_implicit_diff} shows that the gradient of the dual objective \emph{depends only on explicit derivatives of $g_\theta$}, evaluated at samples from $\mu$ and at the proximal points $y_\theta^\star(z)$. In particular, \emph{no implicit differentiation} (computation of $\frac{dy^\star}{d\theta}$) through the inner optimization problem is required. The cancellation follows directly from first-order optimality of the proximal operator and can be interpreted as an instance of the envelope theorem. 


\newcommand{\dgytilde}{\frac{\partial g(\tilde{y}_{\theta}(z))}{\partial \theta}}
\newcommand{\dgystar}{\frac{\partial g(y^*_{\theta}(z))}{\partial \theta}}
\newcommand{\dgx}{\frac{\partial g(x_{\theta})}{\partial \theta}}
\newcommand{\ytilde}{\tilde{y}_{\theta}(z)}
\newcommand{\ystar}{y_{\theta}^*(z)}
\newcommand{\xitheta}{\xi_{\theta}(z)}

\subsection{Convergence}
{\color{black}
In this section, we establish that stochastic gradient descent converges to a stationary point of the loss function even when the inner fixed-point problem is solved only approximately up to a prescribed error tolerance.

In the idealized setting where the fixed-point problem is solved exactly, Lemma~\ref{lem:no_implicit_diff} shows that implicit differentiation is not required for computing gradients with respect to $\theta$. However, this assumption does not hold in practice, as the proximal fixed-point problem must be solved using a finite number of iterations, leading to inexact inner solutions.
As a result, the optimization dynamics are governed by approximate proximal evaluations, and analyzing their effect on convergence is essential.

For ease of presentation, the detailed proofs are deferred to Appendix~\ref{appen:proof}.
}
% Thus, provided that the fixed point problem is solved exactly, JFB is not needed to compute the gradient with respect to $\theta$ of the loss function. This does not occur in practice, but it is proved in this section that stochastic gradient descent (SGD) converges to a local minimum of the loss function if the fixed point problem is solved within a specified error tolerance. For ease of presentation, we provide the proofs in Appendix \ref{appen:proof}.

\subsubsection{Approximate Gradient}
Suppose the fixed point problem is not solved exactly, i.e. instead of the true fixed point $y^*$ an approximate fixed point $\tilde{y}$ is computed.
Then, the stochastic gradient $\tilde{d}_{\theta}$ is given by
\begin{equation}
    \tilde{d}_{\theta} = \mathbb{E}_{x \sim \mu}\left[\frac{\partial g(x)}{\partial \theta}\right] - \mathbb{E}_{z \sim \nu}\left[\frac{\partial g(\tilde{y})}{\partial \theta}\right].
    \label{eq:d_tilde}
\end{equation}

\subsubsection{Preliminary Assumptions and Fixed Point Computation Lemma}

\begin{assumption}
The objective function $L(\theta)$ is bounded from below by $L_{inf}$ in some open subset of its domain.. The function $g$ is $C^1$ with respect to all variables and the gradients of $g$ with respect to $\theta$ and $y$, are $L_\theta$- and $L_y$-Lipschitz, respectively. Furthermore, assume the $2^{nd}$ order partial derivatives of $g$ with respect to $\theta$ exist.
\label{assumption:g_lipschitz}
\end{assumption}

From this point onwards, denote $x \sim \mu$ and $z \sim \nu$ as just $x$ and $z$, respectively.

\newcommand{\lemmaFixedPoint}[1]{
Under Assumption ~\ref{assumption:g_lipschitz}, suppose the approximate fixed point $\tilde{y}$ satisfies
\begin{equation}
\|(\tilde{y} - z) + \nabla_{y}g_{\theta}(\tilde{y})\| < \epsilon_p.
\end{equation}
Then, $\exists \epsilon > 0$ such that
\begin{equation}
\|y^* - \tilde{y}\| < \epsilon.
\end{equation}}
\begin{lemma}
\label{lemma:approx_fixed_point}
\lemmaFixedPoint{main}
\end{lemma}

\subsubsection{Upper Bound on $2^{nd}$ Moments of Exact and Approximate Gradients}
\newcommand{\lemmasecondmoment}[1]{
Under the assumptions of Lemma ~\ref{lemma:approx_fixed_point} along with $\int\|x\|^2d\mu(x) < \infty$ and $\int\|z\|^2d\nu(z) < \infty$,
\begin{equation}
\mathbb{E}_x\left[\left\|\frac{\partial g(x_{\theta})}{\partial \theta}\right\|^2\right] < \infty, \ \mathbb{E}_z\left[\left\|\frac{\partial g(\tilde{y_{\theta}}(z))}{\partial \theta}\right\|^2\right] < \infty, \text{ and } \mathbb{E}_z\left[\left\|\frac{\partial g(y_{\theta}^*(z))}{\partial \theta}\right\|^2\right] < \infty.
\end{equation}}

\begin{lemma}
\label{lemma:2nd_moment}
\lemmasecondmoment{main}
\end{lemma}

\newcommand{\theoremnormsquared}[1]{
Under the assumptions of Lemma ~\ref{lemma:2nd_moment} The norm of the true gradient $\nabla_{\theta}L$ and the stochastic gradient $\tilde{d}_{\theta}$ squared is bounded above, i.e. $\exists 0 < \tilde{M} < +\infty$ and $\exists 0 < M_{L} < +\infty$ such that $\forall \theta$
\begin{equation}
\|\tilde{d}_{\theta}\|^2 = \left\|\mathbb{E}_x\left[\dgx\right] - \mathbb{E}_z\left[\dgytilde\right]\right\|^2 \leq \tilde{M}
\end{equation} and
\begin{equation}
\|\nabla_{\theta}L\|^2 = \left\|\mathbb{E}_x\left[\dgx\right] - \mathbb{E}_z\left[\dgytilde\right]\right\|^2 \leq M_L.
\end{equation}}

\begin{theorem}
\label{theorem:norm_squared}
\theoremnormsquared{main}
\end{theorem}

\subsubsection{Approximate Gradient is a Descent Direction}
\newcommand{\theoremlowerboundinnerproduct}[1]{
Under the assumptions of Theorem ~\ref{theorem:norm_squared}, suppose the norm of the difference between the exact and computed fixed point satisfies, $\forall \theta$ and $\forall z, \  \|\tilde{y}_{\theta}(z) - y_{\theta}^*(z)\| < \epsilon \leq \frac{V\|\nabla_{\theta}L\|^2}{L_{\theta}\left(\sqrt{M_x} + \sqrt{M_z}\right)}$, for some $0 < V < 1$. Then, $\exists U = 1-V > 0$ such that $\forall \theta$ 
\begin{equation}
\langle\nabla_{\theta}L,\tilde{d}_{\theta} \rangle \geq U\|\nabla_{\theta}L\|^2.
\end{equation}}

\begin{theorem}
\theoremlowerboundinnerproduct{main}
\label{theorem:lower_bound_on_inner_product}
\end{theorem}

\subsubsection{Convergence Results}
We begin by introducing notation to formalize the stochasticity arising in the training process. Let $\{\xi_j\}_{j\ge0}$ denote a sequence of independent random variables representing the sampling procedure used to construct the stochastic gradient $\tilde{d}_{\xi_j}(\theta)$. In particular, $\xi_j$ corresponds to the random draw of initial conditions $x \sim \rho$ used to compute the stochastic gradient update at iteration $j$.

We analyze the convergence of SGD when $\tilde{d}_{\theta}$ is used as a stochastic gradient surrogate. Specifically, we consider the iterative scheme
\begin{equation}
\theta_{j+1} = \theta_j - \alpha_j \tilde{d}_{\xi_j}(\theta_j), \qquad j \ge 0,
\label{eq:SGD_iteration}
\end{equation}
for minimizing the loss function over $\theta \in \mathbb{R}^p$. Here, $\tilde{d}_{\xi_j}(\theta_j)$ denotes the JFB update computed using either a single sample or a minibatch of samples with corresponding learning rate $\alpha_j$ at iteration $j$.

Following the notation of ~\cite{bottou2018optimization}, we use $\mathbb{E}_{\xi_j}[\cdot]$ to denote the conditional expectation with respect to the randomness at iteration $j$, given the current iterate $\theta_j$. Since $\theta_j$ depends on the sequence of random variables $\{\xi_0, \xi_1, \ldots, \xi_{j-1}\}$, we also consider the total expectation of the objective with respect to all prior randomness, which we write as
\begin{equation}
\mathbb{E}[\mathbb{E}_x[J_x(\theta_j)]]
=
\mathbb{E}_{\xi_0}\Big[
\mathbb{E}_{\xi_1}\big[
\cdots
\mathbb{E}_{\xi_{j-1}}\big[
\mathbb{E}_x[J_x(\theta_j)]
\big]
\cdots
\big]
\Big].
\end{equation}

With this notation in place, we establish the following Lemma, which is used to prove the main result.

\newcommand{\lemmaExpectationDescent}[1]{
Under the assumptions of Theorem ~\ref{theorem:lower_bound_on_inner_product}, the SGD iterations ~\eqref{eq:SGD_iteration} satisfy
\begin{equation}
\mathbb{E}_{\xi_j}[L(\theta_{j+1})] - L(\theta_j) \leq -\alpha_j U \|\nabla_{\theta}L(\theta_j)]\|^2 + \alpha_j^2L_{\theta}\tilde{M}.
\end{equation}
If $0 < \alpha_j \leq \frac{UM_L}{L_{\theta} \tilde{M}}$, then it follows that $\mathbb{E}_{\xi_j}[L(\theta_{j+1})] - L(\theta_j) \leq 0$.}
\begin{lemma}
\label{lemma:expectation_descent}
\lemmaExpectationDescent{main}
\end{lemma}

With this result, the main results of this paper can be proven.

\newcommand{\theoremExpectationCesaroConvergence}[1]{
Suppose the sequence of learning rates $\{\alpha_j\}_{j=0}^{\infty}$ is monotonically decreasing and satisfies $\sum_{j=0}^{\infty} \alpha_j = \infty$, $\sum_{j=0}^{\infty} \alpha_{j}^2 < \infty$, and $0 < \alpha_0 \leq \frac{UM_L}{L_{\theta} \tilde{M}}$. Let $A_K= \sum_{j=0}^{K-1} \alpha_j$. Then, under the assumptions of Lemma ~\ref{lemma:expectation_descent} the SGD iteration \eqref{eq:SGD_iteration} satisfies
\begin{equation*}
\lim_{K \rightarrow \infty} \mathbb{E} \left[ \frac{1}{A_K}\sum_{j=0}^{K} \alpha_j \left\| \nabla_{\theta}L(\theta_j) \right\|^2 \right] = 0.
\end{equation*}}

\begin{theorem}
\label{theorem:expectation_cesaro_convergence}
\theoremExpectationCesaroConvergence{main}
\end{theorem}

In other words, the weighted Cesaro sum of the sequence $\{ \left\| \nabla_{\theta}L(\theta_j)\right\|^2\}_{j=0}^{\infty}$ converges in (total) expectation to 0. 
Using Theorem~\ref{theorem:expectation_cesaro_convergence}, one can then use standard SGD analysis to show the following theorem and corollary.

\newcommand{\theoremExpectationLiminf}[1]{
Under the assumptions of Theorem ~\ref{theorem:expectation_cesaro_convergence}, the SGD iteration \eqref{eq:SGD_iteration} satisfies
\begin{equation*}
\liminf_{j \rightarrow \infty} \mathbb{E}\left[\left\| \nabla_{\theta}L(\theta_j)\right\|^2 \right] = 0.
\end{equation*}
}

\begin{theorem}
\label{theorem:expectation_liminf}
\theoremExpectationLiminf{main}
\end{theorem}
Using Theorem ~\ref{theorem:expectation_cesaro_convergence}, we can also prove convergence in probability to a critical point.
\newcommand{\corollaryConvergenceProbability}[1]{
Suppose the assumptions of Theorem ~\ref{theorem:expectation_cesaro_convergence} hold. For any $K \in \mathbb{N}$ let $j(K) \in \{0,1,...,K\}$ represent a random index chosen with probabilities proportional to $\{\alpha_j\}_{j=0}^{K}$. Then, $\{\left\| \nabla_{\theta}L(\theta_j)\right\|\}_{j=0}^{K} \rightarrow 0$ as $K \rightarrow \infty$ in probability.}
\begin{corollary}
\label{corollary:convergence_probability}
\corollaryConvergenceProbability{main}
\end{corollary}

% \section{Preliminary: Fixed-Point Framework for HJ PDEs}

% We consider the Hamilton--Jacobi (HJ) equation
% \begin{equation}\label{eq:hj_pre}
% \begin{cases}
% u_t + H(\nabla u) = 0, \\
% u(x,0) = g(x),
% \end{cases}
% \end{equation}
% where the Hamiltonian $H:\mathbb{R}^d \to \mathbb{R}$ is convex and the initial condition 
% $g:\mathbb{R}^d \to \mathbb{R}$ is convex and Lipschitz differentiable.

% Under standard assumptions, the unique viscosity solution of \eqref{eq:hj_pre} admits the Hopf--Lax representation
% \begin{equation}\label{eq:hopf_lax_pre}
% u(x,t) = \inf_{y \in \mathbb{R}^d} \left\{ t\, H^\ast\!\left(\frac{x-y}{t}\right) + g(y) \right\},
% \end{equation}
% where $H^\ast$ denotes the Legendre transform of $H$.

% \paragraph{Fixed-point characterization \cite{park2025fixed}}
% The minimizer $y^\ast$ of \eqref{eq:hopf_lax_pre} satisfies the first-order optimality condition
% \begin{equation}\label{eq:fixed_point_pre}
% y^\ast = x - t\, \nabla H^\ast(\nabla g(y^\ast)).
% \end{equation}
% This defines a fixed-point operator
% \begin{equation}
% F(y) := x - t\, \nabla H^\ast(\nabla g(y)),
% \end{equation}
% leading to the Picard iteration
% \begin{equation}\label{eq:fp_iter_pre}
% y^{(k+1)} = F(y^{(k)}), \qquad k = 0,1,2,\dots
% \end{equation}
% which approximates the Hopf--Lax minimizer.  
% This formulation transforms the variational Hopf--Lax problem into a \emph{simple, mesh-free, gradient-free, parallelizable} iteration, enabling efficient evaluation of high-dimensional HJ equations.

% \begin{theorem}[Convergence Guarantee]\label{thm:fixed_pt}
% If $H$ and $g$ are convex and continuously differentiable with Lipschitz gradients,
% and the time step $t$ satisfies
% \begin{equation}
% t\,L_H L_g < 1,
% \end{equation}
% then $F$ is a contraction, and iteration \eqref{eq:fp_iter_pre} converges globally to the unique minimizer of the Hopf--Lax formula.
% \end{theorem}

% \section{OT via Hopf--Lax and Fixed-Point Iteration}

% We consider OT with quadratic cost
% \begin{equation}
% c(x,z) \;=\; \frac12 \|x-z\|^2,
% \end{equation}
% between probability measures $\mu$ and $\nu$ on $\mathbb{R}^d$.
% Throughout this section, we work with a convex potential
% $g:\mathbb{R}^d \to \mathbb{R}$.

% \subsection{Hopf--Lax formulation}

% Consider the Hamilton--Jacobi equation
% \begin{equation}
% \label{eq:HJ}
% \partial_t u(x,t) + \frac12 \|\nabla u(x,t)\|^2 = 0,
% \qquad
% u(x,0) = g(x).
% \end{equation}
% For the quadratic Hamiltonian $H(p)=\tfrac12\|p\|^2$, the Hopf--Lax formula gives
% \begin{equation}
% \label{eq:hopf-lax}
% u(z,t)
% =
% \inf_{x \in \mathbb{R}^d}
% \left\{
% g(x) + \frac{1}{2t}\|z-x\|^2
% \right\}.
% \end{equation}
% In particular, at $t=1$ we obtain the quadratic inf-convolution
% \begin{equation}
% \label{eq:inf-conv}
% u(z,1)
% =
% \inf_{x \in \mathbb{R}^d}
% \left\{
% \frac12\|z-x\|^2 + g(x)
% \right\}.
% \end{equation}

% \subsection{Relation to the $c$-transform}

% For the quadratic cost $c(x,z)=\tfrac12\|x-z\|^2$, the $c$-transform of a function
% $\phi$ is defined by
% \begin{equation}
% \phi^c(z)
% :=
% \inf_x \{ c(x,z) - \phi(x) \}.
% \end{equation}
% Therefore,
% \begin{equation}
% \label{eq:c-transform-minus}
% (-g)^c(z)
% =
% \inf_x \left\{
% \frac12\|x-z\|^2 + g(x)
% \right\}.
% \end{equation}
% Comparing with \eqref{eq:inf-conv}, we see that
% \begin{equation}
% \label{eq:u1-c}
% u(z,1) = (-g)^c(z).
% \end{equation}
% Thus, the Hopf--Lax solution at time $t=1$ coincides with the
% $c$-transform of $-g$.

% \subsection{Implicit characterization of the minimizer}

% Let $y^\star(z)$ denote a minimizer of \eqref{eq:inf-conv}.
% The first-order optimality condition gives
% \begin{equation}
% 0 = -(z-y^\star(z)) + \nabla g(y^\star(z)),
% \end{equation}
% or equivalently,
% \begin{equation}
% \label{eq:fixed-point}
% y^\star(z) = z - \nabla g(y^\star(z)).
% \end{equation}
% This provides an implicit characterization of the minimizer as a fixed point.

% \subsection{Fixed-point iteration}

% Given $z \in \mathbb{R}^d$, we compute $y^\star(z)$ via Picard iteration:
% \begin{equation}
% \label{eq:picard}
% y^{k+1} = z - \nabla g(y^k),
% \qquad
% k = 0,1,2,\dots
% \end{equation}
% Under standard assumptions (e.g.\ $\nabla g$ Lipschitz with constant $L<1$),
% this iteration converges to the unique fixed point $y^\star(z)$.

% \subsection{Evaluation of the envelope}

% Once $y^\star(z)$ is obtained, the value of the envelope
% $(-g)^c(z)=u(z,1)$ can be recovered as
% \begin{equation}
% \label{eq:envelope-value}
% u(z,1) = (-g)^c(z)
% =
% \frac12\|z-y^\star(z)\|^2 + g(y^\star(z)).
% \end{equation}

% \subsection{Transport maps}
% Assuming $g$ is convex and sufficiently regular, the OT map
% from $\mu$ to $\nu$ is given by ~\cite[Theorem 10.28]{villani2008optimal} (\textcolor{red}{we might be able to show this in a few lines)}
% \begin{equation}
% \label{eq:forward-map}
% T_{\mu\to\nu}(x) = x + \nabla g(x),
% \qquad x \sim \mu,
% \end{equation}
% and the inverse map from $\nu$ to $\mu$ can be expressed as
% \begin{equation}
% \label{eq:backward-map}
% T_{\nu\to\mu}(z) = y^\star(z),
% \qquad z \sim \nu,
% \end{equation}
% where 
% \begin{equation}
%     \label{eq:y_star}
%     y^\star(z) = \argmin_y g(y) + \frac{1}{2}\|z-y\|^2.
% \end{equation}
% Note this gives an implicit characterization of $y^\star(z)$, which can be computed using a fixed point iteration based on, e.g., gradient descent. The optimality/fixed point condition is given by 
% \begin{equation}
%     \label{eq:ystar_fixedpt_condition}
%     \begin{split}
%     & 0 = \nabla g(y^\star) - (z - y^\star)
%     \\
%     \implies & y^\star = z - \nabla g(y^\star).
%     \end{split}
% \end{equation}

% \subsection{Dual objective}

% Let $\mu,\nu$ be probability measures on $\mathbb{R}^d$ and let
% \begin{equation}
% c(x,z) := \frac12\|x-z\|^2.
% \end{equation}
% The Kantorovich dual for the quadratic OT cost~\cite[Thm 5.10]{villani2008optimal} is
% \begin{equation}
% \label{eq:W2-dual}
% W_2^2(\mu,\nu)
% =
% \sup_{\phi \ \text{(proper)}}
% \left\{
% \mathbb{E}_{x\sim\mu}[\phi(x)] + \mathbb{E}_{z\sim\nu}[\phi^c(z)]
% \right\},
% \qquad
% \phi^c(z):=\inf_{y\in\mathbb{R}^d}\{c(y,z)-\phi(y)\}.
% \end{equation}

% \paragraph{Choice of parameterization.}
% We parameterize a convex function $g_\theta$ (e.g.\ ICNN) and set
% \begin{equation}
% \phi_\theta(x) := -g_\theta(x).
% \end{equation}
% Then
% \begin{equation}
% \label{eq:ct-of-minusg}
% \phi_\theta^c(z)=(-g_\theta)^c(z)
% =
% \inf_{y\in\mathbb{R}^d}\left\{\frac12\|y-z\|^2 + g_\theta(y)\right\}.
% \end{equation}

% \paragraph{Hard fixed-point constraint (exact minimizer).}
% Let $y_\theta^\star(z)$ denote a minimizer of \eqref{eq:ct-of-minusg}. If $g_\theta$ is
% differentiable, first-order optimality gives the implicit equation
% \begin{equation}
% \label{eq:hard-fp}
% y_\theta^\star(z)
% =
% z-\nabla g_\theta\!\big(y_\theta^\star(z)\big),
% \qquad z\sim \nu \ \text{a.e.}
% \end{equation}
% We will enforce \eqref{eq:hard-fp} as a hard constraint (i.e.\ $y_\theta^\star(z)$ is defined as the
% exact fixed point / exact minimizer).

% \paragraph{Recovering the $c$-transform value from the minimizer.}
% Substituting $y_\theta^\star(z)$ into \eqref{eq:ct-of-minusg} yields
% \begin{equation}
% \label{eq:ct-value}
% (-g_\theta)^c(z)
% =
% \frac12\|y_\theta^\star(z)-z\|^2 + g_\theta\!\big(y_\theta^\star(z)\big).
% \end{equation}

% \paragraph{Training loss.}
% Maximizing the dual \eqref{eq:W2-dual} over $\phi_\theta=-g_\theta$ is equivalent to minimizing the
% negative dual objective
% \begin{equation}
% \label{eq:loss-dual}
% \mathcal{L}(\theta)
% =
% \mathbb{E}_{x\sim\mu}\big[g_\theta(x)\big]
% -
% \mathbb{E}_{z\sim\nu}\left[
% \frac12\|y_\theta^\star(z)-z\|^2
% +
% g_\theta\!\big(y_\theta^\star(z)\big)
% \right],
% \quad \text{s.t. } y_\theta^\star(z)= \argmin_y \frac{1}{2}\|y-z\|^2 + g_\theta(y)
% \end{equation}
% Minimizing $\mathcal{L}(\theta)$ is equivalent to maximizing the quadratic OT dual over the
% parameterized family $\phi_\theta=-g_\theta$, while enforcing dual feasibility through the hard
% fixed-point condition \eqref{eq:hard-fp}.

\subsection{Class Conditional OT}

In many applications, the source and target measures admit class-wise decompositions
\begin{equation}
\mu = \sum_{k=1}^K \pi_k \mu_k,
\qquad
\nu = \sum_{k=1}^K \pi_k \nu_k,
\end{equation}
where \(\mu_k\) and \(\nu_k\) denote the distributions restricted to class \(k\).  
class-conditional OT (CCOT) enforces that mass is transported only within corresponding classes by solving, for each \(k\),
\begin{equation}
T_k
=
\argmin_{T_\# \mu_k = \nu_k}
\mathbb{E}_{x\sim \mu_k} \Big[ \tfrac12 \| x - T(x) \|^2 \Big].
\end{equation}

Rather than defining separate potentials \(\{ g_k \}_{k=1}^K\) or assuming well-separated supports, we parameterize a single class-conditional potential $g_\theta(x, k) : \mathbb{R}^d \times \{0, \dots, K-1\} \to \mathbb{R}$,
where \(k\) is a class label. The network receives the input \(x\) concatenated with a one-hot encoding of \(k\).


Then, the class-conditional dual objective becomes
\begin{equation}
\label{eq:conditional_potential_loss}
\mathcal{L}_{\mathrm{CC}}(\theta)
=
\frac{1}{K} \sum_{k=0}^{K-1} 
\Bigg[
\mathbb{E}_{x \sim \mu_k} \big[ g_\theta(x, k) \big]
-
\mathbb{E}_{z \sim \nu_k} \Big[ \frac12 \| z - y_\theta^\star(z, k) \|^2 + g_\theta(y_\theta^\star(z, k), k) \Big]
\Bigg],
\end{equation}
where the backward OT map for class \(k\) is defined as the fixed-point solution
\begin{equation}
y_\theta^\star(z, k) = \argmin_{y} \Big\{ \frac{1}{2} \| y - z \|^2 + g_\theta(y, k) \Big\}, 
\quad z \sim \nu_k.
\end{equation}

Hence, conditioning the potential on class labels allows a single network to naturally represent class-conditional transport maps while maintaining a unified convex potential.


% \section{Convergence}
% Forward pass is finding the fixed point of the operator
% \begin{equation}
% F_{\theta}(y) = y - \alpha(\nabla g + y - z) 
% \end{equation}
% Applying the implicit function theorem like in the original JFB paper, we arrive at the expression for the true gradient
% \begin{equation}
% \frac{d y^*}{d \theta} = (I - \frac{\partial F_{\theta}}{\partial y})^{-1}\frac{\partial F_{\theta}}{\partial \theta}
% \end{equation}
% The JFB approximation is
% \begin{equation}
% \frac{d y^*}{d \theta} \approx  \frac{\partial F_{\theta}}{\partial
% \theta} 
% \end{equation}
% Assume that we can interchange $\mathbb{E}_{x}[\cdot]$ and $\mathbb{E}_{z}[\cdot]$ for any $x,z$. Then,
% \begin{align}
% \frac{d L}{d \theta} &= \mathbb{E}_{x \sim \mu}[\frac{d g}{d \theta}] - \mathbb{E}_{z \sim \nu}[\frac{d}{d \theta}(\frac{1}{2}\|y^* - z\|^2 + g(y^*))] \\
% &=  \mathbb{E}_{x \sim \mu}[\frac{d g}{d \theta}] - \mathbb{E}_{z \sim \nu}[(y^* - z)\frac{d y^*}{d \theta} + \frac{ \partial g }{\partial y}\frac{d y^*}{d \theta} + \frac{\partial g}{\partial \theta}] \\
% &= \mathbb{E}_{x \sim \mu}\left[\frac{\partial g}{\partial \theta}\right] - \mathbb{E}_{z \sim \nu}\left[\left((y^* - z) + \frac{ \partial g }{\partial y}\right)\frac{d y^*}{d \theta} + \frac{\partial g}{\partial \theta}\right] 
% \end{align}
% The JFB approximation to $\frac{d L}{d \theta}$, denoted $d_{\theta}^{JFB}$, is given by
% \begin{equation}
% d_{\theta}^{JFB} = \mathbb{E}_{x \sim \mu}\left[\frac{\partial g}{\partial \theta}\right] - \mathbb{E}_{z \sim \nu}\left[\left((y^* - z) + \frac{ \partial g }{\partial y}\right)\frac{\partial F_{\theta}}{\partial \theta} + \frac{\partial g}{\partial \theta}\right] 
% \end{equation}

\section{Experiments}\label{sec:experiments}
We evaluate the proposed method across a variety of datasets and experimental settings to assess its performance and efficiency. We compare against existing baseline methods and further conduct ablation studies to analyze the contribution of each component. 

All experiments are implemented in PyTorch and detailed implementation details are provided in Appendix~\ref{appen:implementation_details}. 
The experiments in Section~\ref{sec:mnist_fmnist} are conducted on an NVIDIA RTX Blackwell 6000 GPU, while all remaining experiments are performed on a single NVIDIA TITAN V GPU (12GB).
The implementation of our method is publicly available at:
\url{https://github.com/Yebbi/ImplicitOT}

\subsection{Evaluation on High-dimensional Gaussian Distributions}\label{sec:high_dim_gaussians_rewrite}

Quantitative evaluation of OT methods for general distributions is often challenging due to the lack of closed-form solutions. To address this, we focus on Gaussian distributions, $\mu = \mathcal{N}(\mathbf{0}, \Sigma_\mu)$ and $\nu = \mathcal{N}(\mathbf{0}, \Sigma_\nu)$, for which the OT map admits an analytical solution:
\begin{equation}\label{eq:OT_Gaussian}
T^{\nu*}_{\mu}(x) = \Sigma_\mu^{-\frac12} (\Sigma_\mu^{\frac12} \Sigma_\nu \Sigma_\mu^{\frac12})^{\frac12} \Sigma_\mu^{-\frac12}x.
\end{equation}
Following the protocol of \cite{korotin2021neural}, we consider dimensions $d \in [2, 64]$, constructing $\Sigma_\mu$ and $\Sigma_\nu$ with random orthonormal eigenvectors and eigenvalues whose logarithms are sampled uniformly from $[-2, 2]$.


\paragraph{Baseline Methods} 
We compare our method with representative OT approaches, including NOT \cite{korotin2022neural}, which directly parameterizes the transport map, MM-v1 \cite{taghvaei20192, korotin2021neural} and MM \cite{korotin2021neural}, which adopt min--max formulations with and without convexity constraints, respectively, NCF \cite{park2025neural}, which constructs transport maps via the HJ equation, LS \cite{seguy2017large}, a dual OT solver with entropic regularization, and WGAN-QC \cite{liu2019wasserstein}, which optimizes a Wasserstein objective under a quadratic cost. Except for NOT, all methods use the same network architecture as ours.

% We compare our method with several established OT approaches. These include NOT \cite{korotin2022neural}, which directly parameterizes the transport map and is trained using the weak formulation; MM-v1 \cite{taghvaei20192, korotin2021neural}, a min--max framework based on input-convex neural networks (ICNNs) that alternates between optimizing a potential function and its convex conjugate; and MM:R \cite{korotin2021neural}, which also follows a min--max formulation but does not impose convexity, instead learning separate networks for the forward and backward maps using a negative Wasserstein loss with an additional conjugacy regularization. We additionally include NCF \cite{park2025neural}, which constructs transport maps based on the characteristics of the Hamilton--Jacobi equation. Furthermore, we evaluate LS \cite{seguy2017large}, a dual OT solver with entropic regularization, and WGAN-QC \cite{liu2019wasserstein}, which adopts a WGAN architecture with a quadratic cost. Except for NOT, all methods, including ours, employ a shared network architecture for modeling the potential functions.


\paragraph{Evaluation Metrics} 
Performance is evaluated using the \emph{unexplained variance percentage (UVP)} \cite{korotin2019wasserstein}. 
Given a predicted transport map $\hat{T}:\mu \to \nu$ and the ground-truth OT map $T^\ast$, the UVP is defined as
\begin{equation}
\mathcal{L}^2\text{-UVP}\left(\hat{T}\right)
\coloneqq
100 \, \frac{\left\lVert \hat{T} - T^\ast \right\rVert_{L^2(\mu)}}{\mathrm{Var}(\nu)} \; (\%).
\end{equation}
We further report computational efficiency in terms of training and inference time, peak memory consumption, and storage requirements for bidirectional transport maps. 

\paragraph{Results} 
\begin{table}[t]
\centering
\vspace{-2mm}
\caption{\textit{Quantitative evaluation on Gaussian distributions.} UVP $(\downarrow)$  is measured across different OT methods as the data dimension $d$ increases. }
\label{tab:gaussians_uvp}
\scalebox{0.86}{
\begin{tabular}{lcccccc}
\toprule
Method & $d=2$ & $d=4$ & $d=8$ & $d=16$ & $d=32$ & $d=64$ \\
\midrule
NOT &   77.248 & 125.419 & 114.056 & 176.086 & 182.287 & 196.831\\
WGAN-QC   &   1.596 & 5.897 & 31.0367 &  59.314 & 113.237 & 141.407\\
LS    & 5.806 & 9.781 & 15.963 & 25.232 & 41.445 & 55.360 \\
MM-v1   & 0.161 & 0.172 &  0.173 & 0.210 & 0.374 & 0.415 \\
MM:R & 0.012 & 0.048 & 0.117 & 0.202 & 0.354 & 0.604 \\
NCF & \textbf{0.010} & 0.021 & 0.086 & 0.146 & 0.436 & 0.858\\
\textbf{Ours} & 0.013 & \textbf{0.016} & \textbf{0.046} & \textbf{0.053} &  \textbf{0.054}  & \textbf{0.0822} \\
\bottomrule
\end{tabular}}
\vspace{-1em}
\end{table}

\begin{figure}
    \centering
    \includegraphics[width=\linewidth]{Figures/Gaussian_memory_time_storage.pdf}
    \caption{\textit{Computational comparison.} Training time (s/epoch), peak memory (MB) during training, and memory (MB) for storing bidirectional OT maps are reported across models and dimensions.
    }
    \label{fig:landscape_comparison}
    \vspace{-1em}
\end{figure}
Table~\ref{tab:gaussians_uvp} reports the UVP scores across different models and dimensionalities, while Figure~\ref{fig:landscape_comparison} summarizes the corresponding computational costs. 
Compared to all baselines, our method consistently learns substantially more accurate OT maps across all tested dimensions. 
While baseline approaches exhibit a noticeable degradation in accuracy as the dimensionality increases, our model maintains a stable error profile even in high-dimensional settings. 
Notably, a single fixed experimental configuration was used for our model across all dimensions, without any dimension-specific tuning. 
These results highlight the accuracy, robustness, and scalability of our approach for high-dimensional OT.


\subsection{Real-World Physics Data Distributions}\label{sec:uci_physics}

To evaluate the extent to which the proposed method can handle complex and practically relevant distributions, we consider real-world data sets with diverse and heterogeneous characteristics.

\paragraph{Data}
We evaluate on UCI benchmark data sets \cite{uci_repo}, which exhibit diverse dimensions and statistical characteristics, including high-energy physics experiments, household power consumption, and chemical sensor measurements of gas mixtures. Following prior work, we define an OT problem from a standard Gaussian distribution to each empirical data distribution, requiring the model to learn a transport map without paired samples.

% We conduct experiments on several benchmark data sets from the University of California Irvine (UCI) machine learning data repository \cite{uci_repo}, including POWER, GAS, HEPMASS, and MINIBOONE, as well as the BSDS300 data set consisting of natural image patches. These data sets arise from a variety of domains, including high-energy physics experiments, household power consumption, and chemical sensor measurements of gas mixtures. As a result, they exhibit a wide range of statistical properties, dimensionalities, and structural features.

% In these settings, the data are not provided as paired samples from two distributions. Instead, we formulate an OT problem by taking a standard Gaussian distribution as the source and the empirical data distribution as the target. The model therefore learns a transport map that pushes forward samples from the Gaussian reference distribution to the target data distribution. This setup provides a practical test of whether the learned transport map can adapt to heterogeneous and high-dimensional distributions in the absence of ground-truth correspondences.


\paragraph{Baseline Methods} 
We compare against NCF \cite{park2025neural}, as well as continuous-flow models including FFJORD \cite{grathwohl2018ffjord}, RNODE \cite{finlay2020train}, and OT-Flow \cite{onken2021ot}. These methods provide strong baselines for evaluating both transport quality and computational efficiency.

% We compare against NCF, the most recent model that achieves the strongest performance in the Gaussian setting among OT-based approaches. In addition, we include continuous flow-based methods, namely FFJORD \cite{grathwohl2018ffjord}, RNODE \cite{finlay2020train}, and OT-Flow \cite{onken2021ot}, as representative approaches for learning continuous transport dynamics under likelihood-based or OT-inspired objectives. These methods provide strong and standard baselines for evaluating both the quality of learned transport maps and their computational efficiency.
% Our study focuses on OT formulations for distribution matching rather than general-purpose generative modeling. Accordingly, we restrict our comparisons to methods operating within a comparable transport or continuous-flow framework.

\begin{table}[t]
\vspace{-1em}
\centering
\small
\setlength{\tabcolsep}{6pt}
% \renewcommand{\arraystretch}{1.2}
\resizebox{0.95\textwidth}{!}{%
\begin{tabular}{l|c|ccccc}
\hline
Dataset & Metric & FFJORD & RNODE & OT-Flow & NCF & Ours \\
\hline

\multirow{2}{*}{Power ($d=6$)} 
& MMD 
& $4.34\!\times\!10^{-5}$ 
& $5.64\!\times\!10^{-5}$ 
& $4.68\!\times\!10^{-5}$ 
& $2.56\!\times\!10^{-4}$
& $\mathbf{1.74\!\times\!10^{-5}}$ \\
& \#Params 
& 43K & 43K & 18K & 8.25K & 8.25K \\
\hline

\multirow{2}{*}{Gas ($d=8$)} 
& MMD 
& $1.02\!\times\!10^{-4}$ 
& $8.03\!\times\!10^{-5}$ 
& $2.47\!\times\!10^{-4}$ 
& $2.24\!\times\!10^{-4}$
& $\mathbf{7.34\!\times\!10^{-5}}$ \\
& \#Params 
& 279K & 279K & 127K & 8.90K & 8.90K \\
\hline

\multirow{2}{*}{HEPMASS ($d=21$)} 
& MMD 
& $1.58\!\times\!10^{-5}$ 
& $1.58\!\times\!10^{-5}$ 
& $1.58\!\times\!10^{-5}$ 
& $6.84\!\times\!10^{-5}$
& $1.79\!\times\!10^{-5}$ \\
& \#Params 
& 547K & 547K & 72K & 13.06K & 13.06K \\
\hline

\multirow{2}{*}{MINIBOONE ($d=43$)} 
& MMD 
& \textbf{$2.84\!\times\!10^{-4}$}
& \textbf{$2.84\!\times\!10^{-4}$} 
& \textbf{$2.84\!\times\!10^{-4}$} 
& $2.78\!\times\!10^{-3}$
& $1.19\!\times\!10^{-3}$ \\
& \#Params 
& 821K & 821K & 78K & 29.84K & 29.84K \\
\hline

\multirow{2}{*}{BSDS300 ($d=63$)} 
& MMD 
& $6.52\!\times\!10^{-3}$ 
& $1.64\!\times\!10^{-2}$ 
& $4.24\!\times\!10^{-4}$ 
& $6.28\!\times\!10^{-4}$
& $\mathbf{3.30\!\times\!10^{-4}}$ \\
& \#Params 
& 6.7M & 6.7M & 297K & 156K & 156K \\
\hline

\end{tabular}}
\caption{Comparison of distribution matching performance (MMD) and model size (\#parameters) across real-world datasets. Lower MMD indicates better distribution matching.}
\label{tab:uci_physics}
\vspace{-2em}
\end{table}


\paragraph{Evaluation Metrics} 
For evaluation, we adopt the same setup as \cite{onken2021ot}, using the unbiased \emph{Maximum Mean Discrepancy (MMD)} estimator with Gaussian kernels to measure the distance between generated and target distributions. Specifically, we use a kernel of the form
\[
k(x,x') = \sum_{j=1}^K \exp(-\alpha_j \|x - x'\|^2),
\]
where the bandwidth parameters $\{\alpha_j\}$ are selected as in prior work. Given two sets of samples $\{x_i\}_{i=1}^{n_1}$ and $\{y_j\}_{j=1}^{n_2}$, the estimator is given by
\[
\mathrm{MMD}^2 = \frac{1}{n_1(n_1-1)} \sum_{i \neq j} k(x_i,x_j)
+ \frac{1}{n_2(n_2-1)} \sum_{i \neq j} k(y_i,y_j)
- \frac{2}{n_1 n_2} \sum_{i,j} k(x_i,y_j).
\]

To assess computational efficiency, we additionally measure the number of trainable parameters in each model.
% \paragraph{Results}
% The results are summarized in Table~\ref{tab:uci_physics}, while quantitative comparisons for Gas, Hepmass, and Miniboone are provided in Figures~\ref{fig:gas}, \ref{fig:hepmass}, and \ref{fig:miniboone}. 
% Since these datasets are high-dimensional, we visualize two-dimensional slices along informative coordinate axes to better illustrate the various structure of the distributions.

% As shown in the figures, the target distributions exhibit highly complex and heterogeneous structures, yet the proposed method is able to accurately capture their geometric characteristics across different datasets. 
% Despite using substantially smaller network architectures, our method achieves on-par or better performance in several benchmarks against continuous flow-based models. While we observe slightly higher MMD in some settings, this comparison should be interpreted in light of the fact that our models are significantly more compact. Moreover, continuous flow-based approaches are primarily developed within likelihood-based generative modeling frameworks and are not directly designed to address general OT problems between arbitrary distributions. In addition, they rely on solving continuous-time dynamics via numerical ODE integration at inference time, which introduces additional computational overhead.
% In contrast, our formulation directly targets the OT problem in a more general setting and does not require iterative ODE solving during sampling, resulting in significantly more efficient inference. Furthermore, compared to NCF, a representative OT-based method, our approach demonstrates consistently better performance on complex, high-dimensional distributions, indicating improved robustness in challenging non-Gaussian settings.
% Overall, these results highlight the generality, accuracy, and computational efficiency of the proposed method.

\paragraph{Results}
Results are summarized in Table~\ref{tab:uci_physics}, with additional visualizations for GAS, HEPMASS, and MINIBOONE in Figures~\ref{fig:gas}, \ref{fig:hepmass}, and \ref{fig:miniboone}. Since the datasets are high-dimensional, we visualize two-dimensional projections along informative coordinate axes.

The proposed method accurately captures the structure of complex target distributions and achieves comparable or better performance than continuous flow-based baselines, despite using smaller network architectures. In contrast to flow-based models, which require continuous-time ODE integration at inference, our method directly solves the OT problem and enables more efficient sampling. Compared to NCF, our approach shows consistently stronger performance on challenging high-dimensional distributions, indicating improved scalability.


\begin{figure}[t]
\centering

\begin{tikzpicture}

% ======================
% spacing control
% ======================
\def\xgap{3.75}   % column spacing 
\def\ygap{-3.2}  % row spacing 

\def\imgw{0.24\linewidth} % image size

% ======================
% columns
% ======================
\foreach \j/\col in {1/0,2/1,4/2,6/3} {

    % ----------------------
    % TRUE row
    % ----------------------
    \node[anchor=center] (T\col)
    at (\col*\xgap, 0)
    {\includegraphics[width=\imgw]{Figures/gas/5_\j_real.png}};

    % ----------------------
    % PRED row
    % ----------------------
    \node[anchor=center] (P\col)
    at (\col*\xgap, \ygap)
    {\includegraphics[width=\imgw]{Figures/gas/5_\j_pred.png}};


    % ----------------------
    % bottom label
    % ----------------------
    \node[align=center] at (\col*\xgap, -5.2)
    {$(5,\j)$-slice};
}

% ======================
% row labels
% ======================
\node[rotate=90, align=center]
at (-2.1, 0.0)
{True};

\node[rotate=90, align=center]
at (-2.1, -3.2)
{Predicted};

\end{tikzpicture}

\caption{
Each column shows a different slice of the UCI Physics Gas dataset, comparing real (top) and predicted (bottom) samples.
}\label{fig:gas}
\vspace{-1em}
\end{figure}


\begin{figure}[t]
\centering

\begin{tikzpicture}

% ======================
% spacing control
% ======================
\def\xgap{3.75}   % column spacing 
\def\ygap{-3.2}  % row spacing 

\def\imgw{0.24\linewidth} % image size

% ======================
% columns
% ======================
\foreach \j/\col in {3/0,6/1,18/2,19/3} {

    % ----------------------
    % TRUE row
    % ----------------------
    \node[anchor=center] (T\col)
    at (\col*\xgap, 0)
    {\includegraphics[width=\imgw]{Figures/hepmass/\j_21_real.png}};

    % ----------------------
    % PRED row
    % ----------------------
    \node[anchor=center] (P\col)
    at (\col*\xgap, \ygap)
    {\includegraphics[width=\imgw]{Figures/hepmass/\j_21_pred.png}};


    % ----------------------
    % bottom label
    % ----------------------
    \node[align=center] at (\col*\xgap, -5.2)
    {$(\j, 21)$-slice};
}

% ======================
% row labels
% ======================
\node[rotate=90, align=center]
at (-2.1, 0.0)
{True};

\node[rotate=90, align=center]
at (-2.1, -3.2)
{Predicted};

\end{tikzpicture}

\caption{
Each column shows a different slice of the UCI Physics Hepmass dataset, comparing real (top) and predicted (bottom) samples.
}\label{fig:hepmass}
\vspace{-1em}
\end{figure}

\begin{figure}[t]
    \centering
    \includegraphics[width=1.0\textwidth]{Figures/CCOT/fmnist_to_mnist.pdf}
    % \caption{Sample image generation using our conditional OT map, trained to transport the distribution of Fashion-MNIST images to that of MNIST. The top row shows input images from Fashion-MNIST, while the bottom row displays the corresponding outputs generated by the model in the MNIST domain. Each column is conditioned on the integer label indicated above the top row.}
    \caption{Conditional OT results for Fashion-MNIST $\rightarrow$ MNIST. Top: input images; bottom: transported outputs. Each column corresponds to a class label.}
    \label{fig:fmnist_to_mnist}
    \vspace{-1em}
\end{figure}


\begin{figure}[t]
\centering

\begin{tikzpicture}

% ======================
% spacing control
% ======================
\def\xgap{3.75}   % column spacing 
\def\ygap{-3.2}  % row spacing 

\def\imgw{0.24\linewidth} % image size

% ======================
% columns
% ======================
\foreach \j/\col in {6/0,13/1,15/2,40/3} {

    % ----------------------
    % TRUE row
    % ----------------------
    \node[anchor=center] (T\col)
    at (\col*\xgap, 0)
    {\includegraphics[width=\imgw]{Figures/Miniboone/0_\j_real.png}};

    % ----------------------
    % PRED row
    % ----------------------
    \node[anchor=center] (P\col)
    at (\col*\xgap, \ygap)
    {\includegraphics[width=\imgw]{Figures/Miniboone/0_\j_pred.png}};


    % ----------------------
    % bottom label
    % ----------------------
    \pgfmathtruncatemacro{\jj}{\j + 1}
    \node[align=center] at (\col*\xgap, -5.2)
    {$(1,\jj)$-slice};
}

% ======================
% row labels
% ======================
\node[rotate=90, align=center]
at (-2.1, 0.0)
{True};

\node[rotate=90, align=center]
at (-2.1, -3.2)
{Predicted};

\end{tikzpicture}

\caption{
Each column shows a different slice of the UCI Physics Miniboone dataset, comparing real (top) and predicted (bottom) samples.
}\label{fig:miniboone}
\vspace{-1em}
\end{figure}



\subsection{2D Class-Conditional Transport}\label{sec:ccot_2d}
We present experimental results on a two-dimensional (2D) synthetic dataset consisting of class-labeled samples, designed to evaluate CCOT. 

\paragraph{Data} The dataset is constructed as a mixture of Gaussian components, where each component corresponds to a distinct class label. This setting allows us to assess whether the learned transport respects both inter-class alignment and intra-class structure preservation.

\paragraph{Results} Figure~\ref{fig:class_2d} illustrates the learned transport behavior across three 2D Gaussian mixture datasets. Each data point is associated with a class label, and transport is performed in a class-conditional manner. The results show that our method successfully aligns corresponding classes between source and target distributions while maintaining clear separation between different classes.

In addition to global distribution matching, the learned transport preserves the local geometry within each class, indicating that the model captures both class-level correspondence and fine-grained structural consistency. This suggests that the proposed formulation can naturally extend to conditional settings without requiring architectural modifications or additional networks.

\begin{figure}
    \centering
    \begin{tikzpicture}[every node/.style={font=\small}]

        % ======================================================
        % Row 1: class_2gaussians
        % ======================================================
        \node[anchor=north west, inner sep=0] (data1)
        at (0, 0) {\includegraphics[width=0.25\linewidth]{Figures/CCOT/class_2gaussians_data.png}};
        \node[anchor=north west, inner sep=0] (fwd1)
        at (0.30\linewidth, 0) {\includegraphics[width=0.25\linewidth]{Figures/CCOT/class_2gaussians_forward.png}};
        \node[anchor=north west, inner sep=0] (bwd1)
        at (0.60\linewidth, 0) {\includegraphics[width=0.25\linewidth]{Figures/CCOT/class_2gaussians_backward.png}};

        % ======================================================
        % Row 2: class_2gaussians2
        % ======================================================
        \node[anchor=north west, inner sep=0] (data2)
        at (0, -4.0) {\includegraphics[width=0.25\linewidth]{Figures/CCOT/class_2gaussians2_data.png}};
        \node[anchor=north west, inner sep=0] (fwd2)
        at (0.30\linewidth, -4.0) {\includegraphics[width=0.25\linewidth]{Figures/CCOT/class_2gaussians2_forward.png}};
        \node[anchor=north west, inner sep=0] (bwd2)
        at (0.60\linewidth, -4.0) {\includegraphics[width=0.25\linewidth]{Figures/CCOT/class_2gaussians2_backward.png}};

        % ======================================================
        % Row 3: class_4gaussians
        % ======================================================
        \node[anchor=north west, inner sep=0] (data3)
        at (0, -8.0) {\includegraphics[width=0.25\linewidth]{Figures/CCOT/class_4gaussians_data.png}};
        \node[anchor=north west, inner sep=0] (fwd3)
        at (0.30\linewidth, -8.0) {\includegraphics[width=0.25\linewidth]{Figures/CCOT/class_4gaussians_forward.png}};
        \node[anchor=north west, inner sep=0] (bwd3)
        at (0.60\linewidth, -8.0) {\includegraphics[width=0.25\linewidth]{Figures/CCOT/class_4gaussians_backward.png}};

        % ======================================================
        % Column titles
        % ======================================================
        \node[anchor=south] at ($(data1.north)+(0,0.3)$) {Data};
        \node[anchor=south] at ($(fwd1.north)+(0,0.3)$) {Forward OT};
        \node[anchor=south] at ($(bwd1.north)+(0,0.3)$) {Backward OT};

        % ======================================================
        % Row labels 
        % ======================================================
        \node[rotate=90, anchor=south] at ($(data1.west)+(-1.0,-0)$)
            {\textsf{Cross Ring}};
        \node[rotate=90, anchor=south] at ($(data2.west)+(-1.0,-0)$)
            {\textsf{Horizontal Swapped}};
        \node[rotate=90, anchor=south] at ($(data3.west)+(-1.0,-0)$)
            {\textsf{Four-Modes}};

    \end{tikzpicture}
    \caption{
    \textit{CCOT on Gaussian mixtures.}
    Each row corresponds to a different class-structured problem.
    The first column shows the empirical data distributions,
    the second column visualizes the learned forward transport map,
    and the third column shows the learned backward transport. Colors indicate different distributions, while marker shapes differentiate individual classes within each distribution.
    }
    \label{fig:class_2d}
    \vspace{-1em}
\end{figure}

\subsection{Class-Conditional Image Translation}\label{sec:mnist_fmnist}
We further evaluate our method on CCOT, where the goal is to transport samples between image distributions while preserving semantic class information. This setting tests whether the learned transport map achieves both distributional alignment and class consistency.

\paragraph{Data}
We use MNIST \cite{lecun1998mnist} and Fashion-MNIST (FMNIST) \cite{xiao2017fashion}, each consisting of 10 semantic classes. Since OT in pixel space is ill-posed for image data \cite{pope2021intrinsic}, we perform transport in a shared VAE latent space of dimension 15 and decode transported samples back to the image domain.

% We use two standard benchmark datasets: MNIST \cite{lecun1998mnist} and Fashion-MNIST (FMNIST) \cite{xiao2017fashion}. MNIST consists of grayscale images of handwritten digits, while FMNIST contains images of fashion products such as clothing and shoes. Both datasets contain 60,000 training images and 10,000 test images, with each dataset composed of 10 classes corresponding to their semantic categories.

% % We consider the OT problem, where samples from Fashion-MNIST (FMNIST) \cite{xiao2017fashion} are mapped to MNIST \cite{lecun1998mnist} while preserving class labels. The transport map is conditioned on class information to enforce alignment between corresponding digit categories.



% Direct application of OT in the pixel space is not well-defined due to the fact that the intrinsic dimension of image datasets is substantially lower than the ambient dimension \cite{pope2021intrinsic}. To address this issue, we perform OT in a latent space learned via a variational autoencoder (VAE). Specifically, both FMNIST and MNIST images are encoded into a shared latent representation, and OT is performed in this latent space. We use a latent dimension of 15, which provides a compact representation while preserving the essential semantic structure of the data. The transported latent codes are subsequently decoded back into the image space using the VAE decoder.

\paragraph{Baseline Methods}
Following prior work \cite{asadulaev2022neural, park2025neural}, we compare against representative image translation and OT methods. These include pixel-level translation models, MUNIT \cite{huang2018multimodal} and AugCycleGAN \cite{almahairi2018augmented, zhu2017unpaired}; OTDD flow \cite{alvarez2021dataset, alvarez2021optimizing}, which preserves class structure through gradient flows; and discrete OT (DOT) based on Sinkhorn optimization with Laplacian regularization \cite{cuturi2013sinkhorn, courty2016optimal}. We also evaluate neural OT methods \cite{korotin2022neural, fan2023neural, asadulaev2022neural} under quadratic (W2), weak quadratic (W2,$\gamma$), and class-conditional (FG) costs, as well as NCF \cite{park2025neural}, which solves class-conditional transport via Hamilton--Jacobi dynamics.
% Following prior works \cite{asadulaev2022neural, park2025neural}, we compare against a wide range of representative baselines spanning generative models and OT methods. Pixel-level adaptation methods include one-to-many translation models such as  MUNIT \cite{huang2018multimodal} and AugCycleGAN \cite{almahairi2018augmented, zhu2017unpaired}. 

% For semi-supervised or label-aware transport, we include OTDD flow \cite{alvarez2021dataset, alvarez2021optimizing}, which leverages gradient flows to preserve class structure during transport. We further evaluate General Discrete OT (DOT), implemented using the Sinkhorn algorithm \cite{cuturi2013sinkhorn} with Laplacian cost regularization \cite{courty2016optimal}

% We also consider neural OT methods \cite{korotin2022neural, fan2023neural, asadulaev2022neural}, evaluated under the quadratic cost $\frac{1}{2}\|x - y\|^2$ (denoted W2) the $\gamma$-weak quadratic cost (W2,$\gamma$ with $\gamma=1$), which allows one-to-many mappings, and class-conditional quadratic cost (FG).

%  Finally, we include HJB characteristic-based NCF \cite{park2025neural}, which solve class-conditional transport via Hamilton--Jacobi dynamics.


\paragraph{Evaluation Metrics}
We evaluate image quality using the \emph{Fréchet Inception Distance (FID)}, which measures the discrepancy between feature distributions of real and generated samples:
[
$\mathrm{FID} = |\mu_r - \mu_g|^2 + \mathrm{Tr}\left(\Sigma_r + \Sigma_g - 2(\Sigma_r \Sigma_g)^{1/2}\right)$.
]

To assess class preservation, we follow \cite{asadulaev2022neural} and report \emph{classification accuracy} using a pretrained ResNet-18 classifier \cite{He_2016_CVPR}. Accuracy is computed as the fraction of transported samples whose predicted label matches the source class label.

% To assess the quality of transported images, we evaluate the \emph{Fréchet Inception Distance (FID)}. FID measures the distance between feature distributions of real and generated samples extracted from a pretrained Inception network. Assuming Gaussian approximations of the feature embeddings with means and covariances $(\mu_r, \Sigma_r)$ and $(\mu_g, \Sigma_g)$ for real and generated data, respectively, FID is defined as:
% \[
% \mathrm{FID} = \|\mu_r - \mu_g\|^2 + \mathrm{Tr}\left(\Sigma_r + \Sigma_g - 2(\Sigma_r \Sigma_g)^{1/2}\right).
% \]
% FID is well-suited for our setting as it captures both sample fidelity and diversity, and is sensitive to distributional shifts across classes, making it a standard metric for evaluating image translation quality in CCOT tasks.

% To further assess the class-preserving property of the learned transport map, we follow the evaluation protocol of \cite{asadulaev2022neural} and measure the \emph{classiciation accuracy} using a pretrained ResNet-18 classifier \cite{He_2016_CVPR}, which achieves over 95\% accuracy on the target domain. For each transported sample $T(x, z)$, we obtain the predicted class label and compare it with the ground-truth label of the source sample $x$. A transport is considered correct if the predicted label of the transported sample $T(x)$ matches the label of $x$, thereby measuring class consistency under the learned mapping.

\paragraph{Results}
FID scores and classification accuracies are summarized in Tables~\ref{tab:fid_scores} and \ref{tab:classification_accuracy}, respectively. Moreover, qualitative results are presented in Figure~\ref{fig:fmnist_to_mnist}. As shown in Table~\ref{tab:fid_scores}, our method achieves the lowest FID among all competing baselines. In particular, it significantly outperforms both pixel-level generative models (AugCycleGAN and MUNIT) and classical OT approaches (W2, W2-$\gamma$, DOT, and OTDD flow). Compared to neural OT methods such as GNOT and NCF, our approach consistently yields improved perceptual quality.
As shown in Table~\ref{tab:classification_accuracy}, our method achieves the highest classification accuracy among all baselines, indicating that it not only aligns marginal distributions but also preserves class-level structure more effectively than competing approaches.

This is further illustrated in Figure~\ref{fig:fmnist_to_mnist}, where the proposed method successfully performs class-consistent transport, aligning samples according to their respective class labels while preserving semantic structure across domains. These results suggest that the proposed single-network minimization framework effectively captures both distributional alignment and class-conditional structure in the latent space, leading to more faithful class-conditional transport compared to adversarial and flow-based neural OT methods.


\begin{table}[t]
\centering
\footnotesize
\setlength{\tabcolsep}{3pt}
\renewcommand{\arraystretch}{0.95}
\caption{Accuracy $\uparrow$ of transported samples.}
\resizebox{\linewidth}{!}{
\begin{tabular}{lcccc|ccc|ccc}
\toprule
& \multicolumn{2}{c|}{Pixel-level} 
& \multicolumn{2}{c|}{Discrete / OTDD}
& \multicolumn{3}{c|}{Neural OT}
& \multicolumn{3}{c}{Ours / Neural Flow} \\
\cmidrule(lr){2-3} \cmidrule(lr){4-5} \cmidrule(lr){6-8} \cmidrule(lr){9-11}
Datasets 
& MUNIT & AugCG 
& OTDD & SinkLpL1 
& W2 & W2,$\gamma$ & FG 
& NCF & FG(no $z$) & \textbf{Ours} \\
\midrule
FMNIST $\rightarrow$ MNIST 
& 8.93 & 12.03 
& 10.28 & 10.67 
& 10.96 & 8.02 & 83.72
& 83.42 & 82.79 & \textbf{96.00} \\
\bottomrule
\end{tabular}
}
\label{tab:classification_accuracy}
\end{table}

\begin{table}[t]
\centering
\footnotesize
\setlength{\tabcolsep}{3pt}
\renewcommand{\arraystretch}{0.95}
\caption{FID $\downarrow$ of transported samples.}
\resizebox{\linewidth}{!}{
\begin{tabular}{lcccc|ccc|ccc}
\toprule
& \multicolumn{2}{c|}{Pixel-level} 
& \multicolumn{2}{c|}{Discrete / OTDD}
& \multicolumn{3}{c|}{Neural OT}
& \multicolumn{3}{c}{Ours / Neural Flow} \\
\cmidrule(lr){2-3} \cmidrule(lr){4-5} \cmidrule(lr){6-8} \cmidrule(lr){9-11}
Datasets 
& MUNIT & AugCG 
& OTDD & SinkLpL1 
& W2 & W2,$\gamma$ & FG 
& NCF & FG(no $z$) & \textbf{Ours} \\
\midrule
FMNIST $\rightarrow$ MNIST 
& 7.91 & 26.35 
& $>$100 & $>$100 
& 7.51 & 7.02 & 5.26 
& 18.27 & 7.14 & \textbf{4.57} \\
\bottomrule
\end{tabular}
}
\label{tab:fid_scores}
\vspace{-1em}
\end{table}


% These results indicate that the proposed implicit formulation yields higher-quality transport maps in the latent space. Compared to neural OT models, our approach consistently improves performance while avoiding instability arising from adversarial training and the additional complexity introduced by regularization-based formulations. Overall, the results suggest that our fized-point-based approach enhances both training stability and sample quality.

% % We also consider CCOT for image generation, mapping samples from the Fashion-MNIST (FMNIST) distribution to MNIST while preserving class labels. The transport map is conditioned on these labels to ensure alignment between corresponding classes.
% To avoid the high dimensionality of pixel space, we first train a variational autoencoder (VAE) to learn a low-dimensional latent representation of the data. We use a latent dimension of 15, which is sufficient to capture the underlying structure and perform transport effectively.
% We then encode both FMNIST and MNIST images into this shared latent space and learn a CCOT map between the latent distributions. Transported samples are decoded back into the image space using the VAE decoder. This approach enables efficient and structured class-conditional transport, as illustrated in Figure~\ref{fig:fmnist_to_mnist}, and achieves competitive FID scores (Table~\ref{tab:fid_scores}). Details on the experimental settings can be found in the appendix.


% \begin{table}[t]
% \centering
% \footnotesize
% \setlength{\tabcolsep}{3pt}
% \renewcommand{\arraystretch}{0.95}
% \caption{FID $\downarrow$ of generated samples.}
% \resizebox{\linewidth}{!}{
% \begin{tabular}{lcccccccccc}
% \toprule
% Datasets 
% & MUNIT & AugCG & OTDD & SinkLpL1 
% & W2 & W2,$\gamma$ 
% & FG (no $z$) & FG 
% & NCF
% & \textbf{Ours} \\
% \midrule
% FMNIST $\rightarrow$ MNIST 
% & 7.91 & 26.35 & $>$100 & $>$100 
% & 7.51 & 7.02 
% & 7.14 & 5.26 
% & 18.27
% & \textbf{4.57} 
% \end{tabular}
% }
% \label{tab:fid_scores}
% \end{table}


\subsection{Ablation Studies}\label{sec:ablation}
\subsubsection{Ablation Study on Fixed-Point Tolerance}
In practical implementations, the proximal subproblem in \eqref{eq:prox_def} is solved using an iterative fixed-point scheme, which is terminated once a prescribed tolerance $\epsilon_p$ is reached. While the theoretical analysis assumes exact computation of $y_\theta^\star(z)$, training is carried out using an approximate solution $\tilde{y}_\theta(z)$. It is therefore important to quantify the effect of the stopping tolerance on both accuracy and computational efficiency.

To this end, we perform an ablation study on the 32-dimensional Gaussian OT task presented in Section \ref{sec:high_dim_gaussians_rewrite}, systematically varying the tolerance $\epsilon_p$. For each configuration, we report the forward and backward transport errors (UVP), along with the average training time per epoch.

\begin{table}[h]
\caption{Sensitivity of transport accuracy and computational cost to the fixed-point stopping tolerance $\epsilon_p$.}
\label{tab:tolerance_ablation}
\centering
\resizebox{\columnwidth}{!}{
\begin{tabular}{c|cccccccc}
\hline
 & $10^{-4}$ & $5\times10^{-4}$ & $10^{-3}$ & $5\times10^{-3}$ & $10^{-2}$ & $5\times10^{-2}$ & $10^{-1}$ & $5\times10^{-1}$ \\
\hline
Forward UVP
& $4.909\times10^{-2}$
& $5.125\times10^{-2}$
& $5.435\times10^{-2}$
& $5.492\times10^{-2}$
& $5.449\times10^{-2}$
& $5.622\times10^{-2}$
& $6.255\times10^{-2}$
& $2.229\times10^{-1}$ \\
Backward UVP
& $4.968\times10^{-2}$
& $5.906\times10^{-2}$
& $5.777\times10^{-2}$
& $5.861\times10^{-2}$
& $5.925\times10^{-2}$
& $6.051\times10^{-2}$
& $6.557\times10^{-2}$
& $1.724\times10^{-1}$ \\
Time (s)
& $0.986$
& $0.917$
& $0.856$
& $0.739$
& $0.640$
& $0.516$
& $0.479$
& $0.277$ \\
\hline
\end{tabular}}
\vspace{-1em}
\end{table}

% \begin{table}[h]
% \centering
% \caption{Sensitivity of transport accuracy and computational cost to the fixed-point stopping tolerance.}
% \label{tab:tolerance_ablation}
% \begin{tabular}{c|c|c|c}
% \hline
% Tolerance $\epsilon_p$ & Forward UVP & Backward UVP & Time (s) \\
% \hline
% $1\times10^{-4}$ & $4.909\times 10^{-2}$ & $4.968\times 10^{-2}$ & $0.986$ \\
% $5\times10^{-4}$ & $5.125\times 10^{-2}$ & $5.906\times 10^{-2}$ & $0.917$ \\
% $1\times10^{-3}$ & $5.435\times 10^{-2}$ & $5.777\times 10^{-2}$ & $0.856$ \\
% $5\times 10^{-3}$ & $5.492\times 10^{-2}$ & $5.861\times 10^{-2}$ & $0.739$ \\
% $1\times10^{-2}$ & $5.449\times 10^{-2}$ & $5.925\times 10^{-2}$ & $0.640$ \\
% $5\times 10^{-2}$ & $5.622\times 10^{-2}$ & $6.051\times 10^{-2}$ & $0.516$ \\
% $1\times10^{-1}$ & $6.255\times 10^{-2}$ & $6.557\times 10^{-2}$ & $0.479$ \\
% $5\times 10^{-1}$ & $2.229\times 10^{-1}$ & $1.724\times 10^{-1}$ & $0.277$ \\
% \hline
% \end{tabular}
% \end{table}

The results exhibit a consistent trade-off between computational cost and estimation accuracy, while simultaneously demonstrating that the method remains robust with respect to the choice of fixed-point tolerance. As predicted by the theory, increasing the tolerance $\epsilon_p$ introduces a discrepancy between the approximate solution $\tilde{y}_\theta(z)$ and the exact minimizer $y_\theta^\star(z)$, which leads to some degradation in accuracy. At the same time, a larger tolerance reduces the number of iterations required for convergence of the fixed-point procedure, thereby improving computational efficiency.

Notably, across the range $\epsilon_p \in [10^{-5}, 10^{-2}]$, the degradation in transport accuracy is marginal, indicating that the method is not highly sensitive to moderate inexactness in the proximal computation. This suggests that the learned transport maps are relatively insensitive to approximation errors in the fixed-point solution within this regime. This observation is consistent with the analysis in Section~\ref{lemma:approx_fixed_point}, which guarantees that sufficiently accurate approximate fixed points remain close to the true minimizer.

From a practical perspective, this empirical behavior indicates that, within a reasonable tolerance regime, the approximate fixed-point $\tilde{y}_\theta(z)$ provides a sufficiently accurate surrogate for $y_\theta^\star(z)$ and does not significantly distort the gradient signal used during training. In contrast, when the tolerance becomes excessively large (e.g., $\epsilon_p \geq 10^{-1}$), the approximation error becomes non-negligible and leads to a clear deterioration in performance.

Overall, these findings suggest that solving the inner optimization problem to very high precision is unnecessary in practice. Instead, moderately accurate fixed-point solutions achieve a favorable balance, providing substantial computational savings while maintaining transport accuracy.

\subsubsection{Ablation Study on Potential Parameterization}
To study the effect of potential parameterization, we compare two ICNN variants: (i) a \textbf{Convex Network}, where convexity is enforced via weight projection, and (ii) a \textbf{Nonconvex Network}, which uses the same architecture without convexity constraints. We further evaluate three activation functions: Tanh, SoftPlus, and CeLU. All experiments are conducted in the 32-dimensional setting. Results are reported in Table~\ref{tab:activation_ablation}, and the fixed-point iteration depth and residual
$|y^{k+1}-z+\nabla g(y^{k+1})|_\infty$
are shown in Figure~\ref{fig:network_ablation}.

% To investigate how different neural parameterizations of the Kantorovich potential $g$ affect the behavior of our method, we conduct an ablation study focusing on architectural convexity. Specifically, we consider the following variants:
% \begin{itemize}
%     % \item \textbf{MLP}: an unconstrained scalar-valued MLP with SoftPlus activation.
%     \item \textbf{Convex Network}: an input-convex neural network (ICNN) where convexity is enforced via weight projection at each iteration.
    
%     \item \textbf{Nonconvex Network}: an ICNN trained without enforcing convexity constraints. This setting preserves the same underlying architecture, allowing us to isolate the effect of convexity enforcement without confounding architectural differences.
% \end{itemize}
% In addition, to examine the effect of nonlinearities, we evaluate each architecture using three different activation functions: Tanh, SoftPlus, and CeLU. This allows us to disentangle the impact of convexity constraints from that of the activation choice on optimization stability and transport performance.

% \begin{table}[t]
% \centering
% \caption{Comparison of network parameterization across dimensions (32D and 64D). We report forward and backward UVP values.}
% \label{tab:ablation_network}
% \begin{tabular}{lcccc}
% \hline
% \multirow{2}{*}{Method} 
% & \multicolumn{2}{c}{$d=32$} 
% & \multicolumn{2}{c}{$d=64$} \\
% \cline{2-5}
% & Forward & Backward & Forward & Backward \\
% \hline
% MLP 
% & $7.839 \times 10^{-1}$ & $8.912 \times 10^{-1}$ 
% & $6.698 \times 10^{-1}$ & $7.346 \times 10^{-1}$ \\
% ICNN w/ Convexity 
% & $6.826  \times 10^{0}$  & $5.794\times 10^{0}$ 
% & $3.161\times 10^{0}$ & $3.397\times 10^{0}$ \\
% ICNN wo/ Convexity
% & $5.4 \times 10^{-2}$ & $5.67 \times 10^{-2}$ 
% & $8.22 \times 10^{-2}$ & $7.73 \times 10^{-2}$ \\
% \hline
% \end{tabular}
% \end{table}
\begin{table}[t]
\centering
\footnotesize
\setlength{\tabcolsep}{4pt}
\renewcommand{\arraystretch}{0.8}
\caption{Comparison of network parameterization.}\label{tab:ablation_network}
\resizebox{0.6\columnwidth}{!}{
\begin{tabular}{lcc|cc}
\toprule
 & \multicolumn{2}{c}{\textbf{Convex}} 
 & \multicolumn{2}{c}{\textbf{Non-convex}} \\
\cmidrule(lr){2-3} \cmidrule(lr){4-5}
\textbf{Activation} 
& Fwd UVP & Bwd UVP 
& Fwd UVP & Bwd UVP \\
\midrule
Tanh 
& $7.173 \times 10^{-1}$ & $7.565 \times 10^{-1}$ 
& $1.730 \times 10^{-1}$ & $2.184 \times 10^{-1}$ \\

SoftPlus 
& $6.857 \times 10^{0}$ & $5.945 \times 10^{0}$ 
& $6.996 \times 10^{-2}$ & $7.551 \times 10^{-2}$ \\

CeLU 
& $6.826 \times 10^{0}$ & $5.794 \times 10^{0}$ 
& $5.4 \times 10^{-2}$ & $5.67 \times 10^{-2}$ \\
\bottomrule
\end{tabular}}
\label{tab:activation_ablation}
\end{table}
% \begin{table}[t]
% \centering
% \caption{Comparison of network parameterization across dimensions (32D and 64D). We report forward and backward UVP values.}
% \label{tab:ablation_network}
% \begin{tabular}{lcccc}
% \hline
% \multirow{2}{*}{Method} 
% & \multicolumn{2}{c}{$d=32$} 
% & \multicolumn{2}{c}{$d=64$} \\
% \cline{2-5}
% & Forward & Backward & Forward & Backward \\
% \hline
% Tanh  w/ Convexity 
% & $7.839 \times 10^{-1}$ & $8.912 \times 10^{-1}$ 
% & $6.698 \times 10^{-1}$ & $7.346 \times 10^{-1}$ \\
% Tanh  wo/ Convexity 
% & $7.839 \times 10^{-1}$ & $8.912 \times 10^{-1}$ 
% & $6.698 \times 10^{-1}$ & $7.346 \times 10^{-1}$ \\
% SoftPlus  w/ Convexity 
% & $7.839 \times 10^{-1}$ & $8.912 \times 10^{-1}$ 
% & $6.698 \times 10^{-1}$ & $7.346 \times 10^{-1}$ \\
% SoftPlus  wo/ Convexity 
% & $7.839 \times 10^{-1}$ & $8.912 \times 10^{-1}$ 
% & $6.698 \times 10^{-1}$ & $7.346 \times 10^{-1}$ \\
% CeLU w/ Convexity 
% & $6.826  \times 10^{0}$  & $5.794\times 10^{0}$ 
% & $3.161\times 10^{0}$ & $3.397\times 10^{0}$ \\
% ICNN wo/ Convexity
% & $5.4 \times 10^{-2}$ & $5.67 \times 10^{-2}$ 
% & $8.22 \times 10^{-2}$ & $7.73 \times 10^{-2}$ \\
% \hline
% \end{tabular}
% \end{table}

% \begin{figure}[t]
%     \centering

%     \begin{subfigure}[t]{0.32\textwidth}
%         \centering
%         \includegraphics[width=\textwidth]{Figures/Ablation/Network_Ablation_UVP.png}
%         \caption{Forward UVP}
%         \label{fig:uvp}
%     \end{subfigure}
%     \hfill
%     \begin{subfigure}[t]{0.32\textwidth}
%         \centering
%         \includegraphics[width=\textwidth]{Figures/Ablation/Network_Ablation_Iterations.png}
%         \caption{Fixed-point iterations}
%         \label{fig:iter}
%     \end{subfigure}
%     \hfill
%     \begin{subfigure}[t]{0.32\textwidth}
%         \centering
%         \includegraphics[width=\textwidth]{Figures/Ablation/Network_Ablation_Residual.png}
%         \caption{Fixed-point residual}
%         \label{fig:res}
%     \end{subfigure}

%     \caption{Ablation study on network parameterizations. We compare MLP (ImplicitNet), ICNN with convexity, and ICNN without explicit convexification across (a) OT map accuracy (UVP), (b) the number of fixed-point iteration to converge, and (c) convergence residual.}
%     \label{fig:network_ablation}
% \end{figure}
\begin{figure}[t]
    \centering
    \includegraphics[width=0.99\linewidth]{Figures/Ablation/Network_Ablation.png}
    \caption{Ablation study on network parameterizations. We compare convex and nonconvex network with different activation functions across (a) OT map accuracy (UVP), (b) the number of fixed-point iteration to converge, and (c) convergence residual.}
    \label{fig:network_ablation}
    \vspace{-1em}
\end{figure}
% We evaluate each model under 32-dimensional case.
% The results are summarized in Table~\ref{tab:ablation_network}, 
% while the evolution of the fixed-point iteration depth and residual
% $ \left\| y^{k+1} - z + \nabla g(y^{k+1}) \right\|_{\infty}$ are illustrated in Figure~\ref{fig:network_ablation}.

\paragraph{Trade-off Perspective.}
Table~\ref{tab:activation_ablation} and Figure~\ref{fig:network_ablation} reveal a trade-off between the expressiveness of the learned potential and the stability of the induced fixed-point iterations. Convexified ICNNs yield the most stable and fastest convergence due to their monotone gradient fields, but their constrained parameterization limits flexibility and leads to inferior OT objective values. While convexity is theoretically well-motivated by the Kantorovich formulation, strong convexification restricts the effective parameter space and introduces optimization bias. 
% As a result, the learned potential lacks the flexibility needed to accurately capture the geometry of the target transport, leading to suboptimal values of the dual objective.

In contrast, nonconvex ICNNs achieve the best overall performance, providing sufficient expressive power while maintaining stable fixed-point dynamics. This suggests that explicit convexification is not necessary in practice and may overly restrict the learned transport geometry.

Activation functions exhibit a similar effect. Tanh consistently performs worse, likely due to gradient saturation, which degrades both optimization and fixed-point convergence. In contrast, SoftPlus and CeLU provide smoother gradient flow and better-conditioned optimization, resulting in improved convergence and transport performance.



% The results in Table~\ref{tab:activation_ablation} and Figure~\ref{fig:network_ablation} highlight a critical trade-off in our method between the expressive power of the neural potential and the stability of the induced fixed-point iterations. Since our approach relies on solving a proximal fixed-point problem at every step, good performance requires both (i) a sufficiently expressive parameterization of the Kantorovich potential g and (ii) stable and efficient convergence of the fixed-point dynamics.

% ICNNs with explicit convexification enforce convexity through architectural constraints, yielding monotone gradient fields and, in principle, well-behaved fixed-point updates. Indeed, we observe that this setting leads to the most stable and fastest convergence in terms of fixed-point iterations. However, such strong enforcement of convexity degrades overall performance. While convexity is theoretically well-motivated by the Kantorovich formulation, strong convexification restricts the effective parameter space and introduces optimization bias. As a result, the learned potential lacks the flexibility needed to accurately capture the geometry of the target transport, leading to suboptimal values of the dual objective.

% In contrast, ICNNs without explicit convexification achieve the best overall performance. In this case, convexity is not explicitly enforced through hard projection. This provides a favorable balance: the model retains sufficient flexibility to approximate complex transport structures while still benefiting from a partially monotone gradient field. Empirically, this leads to maintain stable fixed-point iterations—while achieving substantially better optimization of the OT objective compared to convexified ICNNs.

% The choice of activation function further reveals the importance of this balance. In particular, Tanh consistently underperforms both in terms of objective value and fixed-point convergence. Although Tanh is Lipschitz and bounded, its saturation behavior leads to vanishing gradients, especially in deeper compositions. This results in poorly conditioned gradient fields for $\nabla g_\theta$ which in turn degrades the contraction properties required for stable fixed-point iterations. Consequently, the proximal updates become less reliable, often requiring more iterations and converging to inferior solutions.

% On the other hand, smoother and non-saturating activations such as \textbf{SoftPlus} and \textbf{CeLU} provide better gradient flow and improved conditioning of the optimization landscape. These activations enable richer functional representations while maintaining sufficiently stable dynamics, leading to both improved fixed-point convergence and superior transport performance.

\section{Conclusion}\label{sec:conclusion}
We studied an implicit OT formulation based on a single neural potential, where both forward and backward transport maps are recovered through a unified minimization problem. By exploiting the proximal structure of the dual formulation, the objective reduces to a fixed-point problem whose gradient can be computed without implicit differentiation. We establish stochastic gradient descent convergence under inexact inner fixed-point iterations, which arise in practice due to finite computational budgets.
This yields a simple and stable training procedure that avoids saddle-point optimization, auxiliary networks, and unrolled dynamics, while maintaining strong empirical performance in high-dimensional and class-conditional settings.
% We proposed a novel formulation of quadratic OT based on a single neural potential, in which both forward and backward transport maps are recovered through a unified minimization problem. By exploiting the proximal structure of the dual formulation, the objective reduces to a fixed-point problem whose gradient can be computed without implicit differentiation. This yields a simple and stable training procedure that avoids saddle-point optimization, auxiliary networks, and unrolled dynamics, while maintaining strong empirical performance in high-dimensional and class-conditional settings.

Future work includes extending the evaluation to more diverse and complex datasets, as well as establishing stronger theoretical guarantees, in particular convergence to the true OT map under practical neural network parameterizations. Furthermore, generalizing the proposed framework beyond the quadratic cost to more general cost functions remains an important direction for broadening its applicability.

\appendix
\section{Proofs}\label{appen:proof}

\subsection{Proof of Theorem ~\ref{theorem:optimal_transport_map}}
\begin{proof}
Because $\mu, \nu$ are probability measures on $\mathbb{R}^d$ with finite second moments, and $\mu$ is absolutely continuous with respect to the Lebesgue measure on $\mathbb{R}^d$, all of the assumptions of Brenier's Theorem~\cite[Theorem 2.5.10]{figalli_glaudo_OT} are true. Then by ~\cite[Corollary 2.5.12]{figalli_glaudo_OT} $\exists$ a unique optimal transport map $T: \mathbb{R}^d \rightarrow \mathbb{R}^d$ such that $T_{\#}\mu = \nu$ and $T = \nabla u$ for some convex $u: \mathbb{R}^d \rightarrow \mathbb{R} \cup \{+\infty\}$. \newline \newline
By~\cite[Theorem 10.28]{villani2008optimal}, $T$ satisfies $\nabla g(x) + \nabla_x c(x,T(x)) = 0$ $\mu$-almost everywhere for some c-convex function $g:\mathbb{R}^d \rightarrow \mathbb{R}\cup \{+\infty\}$. Thus, $\mu$-a.e.,
\begin{align}
\nabla g(x) + \nabla_{x}\left(\frac{1}{2}\|x -T(x) \|^2\right) &= 0 \\
\nabla g(x) + x - T(x) &= 0.
\end{align}. 
Thus,
\begin{equation}
    T(x) = x + \nabla g(x).
\end{equation}
Finally, since $u$ and $c(x,y)$ are convex, $g$ must be 1-weakly convex.
% \begin{equation}
% T(x) = \nabla\left(\frac{1}{2}\|x\|^2 + g\right) = x + \nabla g.
% \end{equation}
% Let $u(x) = \frac{1}{2}\|x\|^2 + g(x)$. Because it is the sum of a convex function and a 1-weakly convex function, $u$ is convex. Thus, $T(x) = \nabla u(x)$ and we have already proved that $T$ is the unique optimal transport map from $\mu$ to $\nu$.
\end{proof}

\subsection{Proof of Lemma ~\ref{lemma:approx_fixed_point}}
\begin{proof}
Because $\nabla_yg_{\theta}$ is assumed to be $L_y$ Lipschitz in $y$ it follows that $\nabla_yg_{\theta}$ is continuous in $y$. Therefore, the function
\begin{equation*}
    \tilde{y} \mapsto (\tilde{y} - z) + \nabla_yg_{\theta}(\tilde{y}) 
\end{equation*}
is continuous at all points in its domain. Then, by definition, $\forall \epsilon' > 0$, $\exists \delta' > 0$ such that $\|(\tilde{y} - z) + \nabla_yg_{\theta}(\tilde{y}) - \left( (y^* - z) - \nabla_{y}g_{\theta}(y^*)\right)\| < \epsilon'$ whenever $\|\tilde{y} - y^*\| < \delta'$. Because $y^*$ is an exact minimizer of $\frac{1}{2}\|y-z\|^2 + g_{\theta}(y)$, it follows that whenever $\|\tilde{y} - y^*\| < \delta'$,
\begin{equation*}
\|(\tilde{y} - z) + \nabla_yg_{\theta}(\tilde{y}) - \left( (y^* - z) - \nabla_{y}g_{\theta}(y^*)\right)\| = \|(\tilde{y} - z) + \nabla_yg_{\theta}(\tilde{y}) - 0\| = \|(\tilde{y} - z) + \nabla_yg_{\theta}(\tilde{y})\| < \epsilon'.
\end{equation*}
Thus, the result follows by choosing $\epsilon = \delta'$ corresponding to $\epsilon' = \epsilon_p$.
\end{proof}

\subsection{Proof of Lemma ~\ref{lemma:2nd_moment}}
\begin{proof}
Because $\frac{\partial g}{\partial \theta}$ is $L_{\theta}$-Lipschitz in $\theta$ and by the result of Lemma ~\ref{lemma:approx_fixed_point}, $\exists \epsilon > 0$ such that $\|\tilde{y}_{\theta}(z) - y_{\theta}^*(z)\| < \epsilon$, it follows that
\begin{equation*}
0 \leq \mathbb{E}_z\left[\left\|\frac{\partial g(\tilde{y}_{\theta}(z))}{\partial \theta} - \frac{\partial g(\tilde{y}_{\theta}(z))}{\partial \theta} \right\|^2\right] \leq \mathbb{E}_z \left[ L_g\left\|\tilde{y}_{\theta}(z) - y_{\theta}^*(z)\right\|^2\right] \leq \epsilon L_g.
\end{equation*}
Expanding $\mathbb{E}_z\left[\left\|\frac{\partial g(\tilde{y}_{\theta}(z))}{\partial \theta} - \frac{\partial g(\tilde{y}_{\theta}(z))}{\partial \theta} \right\|^2\right]$, the above inequality becomes
\begin{equation*}
0 \leq \mathbb{E}_z\left[\left\|\frac{\partial g(\tilde{y}_{\theta}(z))}{\partial \theta} \right\|^2\right] -2\mathbb{E}_z\left[\frac{\partial g(\tilde{y}_{\theta}(z))}{\partial \theta}^{\top}\frac{\partial g(y_{\theta}^*(z))}{\partial \theta}\right] + \mathbb{E}_z\left[\left\|\frac{\partial g(\tilde{y}_{\theta}(z))}{\partial \theta}\right\|^2\right] \leq \epsilon L_g.
\end{equation*}
It follows from the above inequality and the assumption that $\int\|x\|^2d\mu(x)$ and $\int\|z\|^2d\nu(z)$ that $\mathbb{E}_z\left[\left\|\frac{\partial g(\tilde{y}_{\theta}(z))}{\partial \theta} \right\|^2\right]$,  $\mathbb{E}_z\left[\frac{\partial g(\tilde{y}_{\theta}(z))}{\partial \theta}^{\top}\frac{\partial g(y_{\theta}^*(z))}{\partial \theta}\right]$, and $\mathbb{E}_z\left[\left\|\frac{\partial g(\tilde{y}_{\theta}(z))}{\partial \theta}\right\|^2\right]$ are all finite. \newline
A near-identical argument shows $\mathbb{E}_x\left[\left\|\frac{\partial g(x_{\theta})}{\partial \theta}\right\|^2\right]$ is finite by picking any two points $x_{\theta,1}, x_{\theta,2}$ such that $\|x_{\theta,1} - x_{\theta,2}\| < \epsilon$ and using the fact that $\frac{\partial g(x_{\theta})}{\partial \theta}$ is $L_g$-Lipschitz.
\end{proof}

\subsection{Proof of Theorem ~\ref{theorem:norm_squared}}
\begin{proof}
Expanding the definition of $\tilde{d}$,
\begin{align}
\|\tilde{d}_{\theta}\|^2 &= \left\|\mathbb{E}_x\left[\dgx\right] - \mathbb{E}_z\left[ \dgytilde \right] \right\|^2 \\
&= \left\| \mathbb{E}_x\left[\dgx \right] \right\|^2 - 2\left\langle \mathbb{E}_x\left[\dgx\right], \mathbb{E}_z\left[ \dgytilde \right] \right\rangle + \left\| \mathbb{E}_z\left[ \dgytilde \right]\right\|^2 \\
&= \mathbb{E}_x\left[\left\| \dgx\right\|^2 \right] - 2\left\langle \mathbb{E}_x\left[ \dgx \right], \mathbb{E}_z\left[ \dgytilde \right] \right\rangle + \mathbb{E}_z\left[ \left\|\dgytilde\right\|^2 \right], 
\label{eq:2nd_moment}
\end{align}
where ~\eqref{eq:2nd_moment} follows by Jensen's inequality. By the result of Lemma ~\ref{lemma:2nd_moment}, $\exists 0 < M_x < +\infty, 0 < M_z < +\infty$ such that $\mathbb{E}_x\left[\left\| \dgx\right\|^2\right] < M_x$ and $\mathbb{E}_z\left[\left\| \dgytilde \right\|^2\right] < M_z$. Applying Jensen's inequality to $\mathbb{E}_x\left[\left\| \dgx\right\|^2\right], \mathbb{E}_z\left[\left\| \dgytilde \right\|^2\right]$ yields
\begin{equation*}
\left\| \mathbb{E}_x\left[\dgx\right] \right\| < \sqrt{M_x} \text{ and } \left\| \mathbb{E}_x\left[\dgytilde\right] \right\| < \sqrt{M_z}.
\end{equation*}
It then follows by the Cauchy-Schwarz inequality that
\begin{equation}
- \left\| \mathbb{E}_x\left[ \dgx \right]\right\| \left\| \mathbb{E}_x\left[\dgytilde\right] \right\| \leq \left\langle \mathbb{E}_x\left[ \dgx \right], \mathbb{E}_z\left[ \dgytilde \right] \right\rangle \leq \left\| \mathbb{E}_x\left[ \dgx \right]\right\| \left\| \mathbb{E}_x\left[\dgytilde\right] \right\|.
\label{eq:cs}
\end{equation}
Combining ~\eqref{eq:2nd_moment} and ~\eqref{eq:cs},
\begin{equation*}
0 \leq \|\tilde{d}_{\theta}\|^2 \leq M_x + M_z + 2\sqrt{M_xM_z} =\tilde{M} < + \infty.
\end{equation*}
An identical argument shows the existence of $M_L$ such that $0 \leq \|\nabla_{\theta}L\|^2 \leq M_L < +\infty$.
\end{proof}

\subsection{Proof of Theorem ~\ref{theorem:lower_bound_on_inner_product}}
\begin{proof}
Expanding the inner product, 
\begin{equation*}
\langle \nabla_{\theta}L, \tilde{d}_{\theta}\rangle = \left\langle \mathbb{E}_x\left[\dgx\right] - \mathbb{E}_z\left[\dgystar\right], \mathbb{E}_x\left[\dgx\right] - \mathbb{E}_z\left[\dgytilde\right]\right\rangle
\end{equation*}
\begin{equation*}
\begin{split}
\langle \nabla_{\theta}L, \tilde{d}_{\theta}\rangle = \left\langle \mathbb{E}_x\left[\dgx\right], \mathbb{E}_x\left[\dgx\right]\right\rangle - \left\langle \mathbb{E}_x\left[\dgx\right], \mathbb{E}_z\left[\dgytilde\right] + \mathbb{E}_z\left[\dgystar\right]\right\rangle + \\ \left\langle\mathbb{E}_z\left[\dgytilde\right],\mathbb{E}_z\left[\dgystar\right]\right\rangle
\end{split}
\end{equation*}

\begin{equation}
\begin{split}
\langle \nabla_{\theta}L, \tilde{d}_{\theta}\rangle = \left\| \mathbb{E}_x\left[\dgx\right]\right\|^2 - \left\langle \mathbb{E}_x\left[\dgx\right], \mathbb{E}_z\left[\dgytilde\right] + \mathbb{E}_z\left[\dgystar\right]\right\rangle + \\\left\langle\mathbb{E}_z\left[\dgytilde\right],\mathbb{E}_z\left[\dgystar\right]\right\rangle.
\label{eq:expand_inner_product}
\end{split}
\end{equation}
$\forall z$, let $\xi_{\theta}(z) = \ytilde - \ystar$ so $\|\xi_{\theta}(z)\| = \epsilon$. Fix $z$. Because the $2^{nd}$ order partial derivatives of $g$ with respect to $\theta$ are assumed to exist, by Taylor's Theorem, $\exists t \in (0,1)$ such that
\begin{equation}
\dgytilde = \dgystar + \nabla_{\theta}^2g(\ystar + t\xi_{\theta}(z))\xi_{\theta}(z).
\end{equation}
Taking expectation with respect to $z$ on both sides, by linearity of expectation,
\begin{equation}
\mathbb{E}_z\left[\dgytilde\right] = \mathbb{E}_z\left[\dgystar\right] + \mathbb{E}_z\left[\nabla_{\theta}^2g(\ystar + t\xi_{\theta}(z))\xi_{\theta}(z)\right].
\end{equation}
Therefore,
\begin{equation*}
    \begin{split}
    \left\langle \mathbb{E}_x\left[\dgx\right],\mathbb{E}_z\left[\dgytilde\right] + \mathbb{E}_z\left[\dgystar\right] \right\rangle = \\
    \left\langle \mathbb{E}_x\left[ \dgx \right],\mathbb{E}_z\left[\dgystar\right] +\mathbb{E}_z\left[\nabla_{\theta}^2g(\ystar + t\xitheta)\xitheta\right] + \mathbb{E}_z\left[\dgystar\right]\right\rangle
    \end{split} 
\end{equation*}

\begin{equation}
    \begin{split}
    \left\langle \mathbb{E}_x\left[\dgx\right], \mathbb{E}_z\left[\dgytilde\right] + \mathbb{E}_z\left[\dgystar\right]\right\rangle= 2\left\langle \mathbb{E}_x\left[\dgx\right], \mathbb{E}_z\left[\dgystar\right]\right\rangle + \\ \left\langle \mathbb{E}_x\left[\dgx\right],\mathbb{E}_z\left[\nabla_{\theta}^2g(\ystar + t\xitheta)\xitheta\right]  \right\rangle
    \end{split}
\label{eq:expand_middle}
\end{equation}
and
\begin{align}
\left\langle \mathbb{E}_z\left[\dgytilde\right], \mathbb{E}_z\left[\dgystar\right]\right\rangle &= \left\langle \mathbb{E}_z\left[\dgystar\right] + \mathbb{E}_z\left[\nabla_{\theta}^2g(\ystar + t\xitheta)\xitheta\right], \mathbb{E}_z\left[\dgystar\right]\right\rangle \\ 
&= \left\| \mathbb{E}_z\left[\dgystar\right]\right\|^2 + \left\langle \mathbb{E}_z\left[\nabla_{\theta}^2g(\ystar + t\xitheta)\xitheta\right], \mathbb{E}_z\left[\dgystar\right]\right\rangle.
\label{eq:expand_right}
\end{align}
Combining ~\eqref{eq:expand_inner_product}, ~\eqref{eq:expand_middle}, and ~\eqref{eq:expand_right} gives
\begin{equation*}
\begin{split}
\langle \nabla_{\theta}L, \tilde{d}_{\theta} \rangle = \left\| \mathbb{E}_x\left[\dgx\right]\right\|^2 - 2\left\langle \mathbb{E}_x\left[\dgx\right],\mathbb{E}_z\left[\dgystar\right]\right\rangle + \left\| \mathbb{E}_z\left[\dgystar\right]\right\|^2 - \\
\left\langle \mathbb{E}_z\left[\nabla_{\theta}^2g(\ystar + t\xi_{\theta}(z))\xitheta\right], \mathbb{E}_x\left[\dgx\right] + \mathbb{E}_z\left[\dgystar\right]  \right\rangle
\end{split}
\end{equation*}
\begin{equation}
\langle \nabla_{\theta}L, \tilde{d}_{\theta} \rangle = \left\| \nabla_{\theta}L\right\|^2 - \left\langle \mathbb{E}_z\left[\nabla_{\theta}^2g(\ystar + t\xi_{\theta}(z))\xitheta\right], \mathbb{E}_x\left[\dgx\right] + \mathbb{E}_z\left[\dgystar\right]  \right\rangle.
\label{eq:complete_the_square}
\end{equation}
By Cauchy-Schwarz and the triangle inequality
\begin{equation*}
\begin{split}
\left\langle \mathbb{E}_z\left[\nabla_{\theta}^2g(\ystar + t\xi_{\theta}(z))\xitheta\right], \mathbb{E}_x\left[\dgx\right] + \mathbb{E}_z\left[\dgystar\right]  \right\rangle \leq \\\left\|\mathbb{E}_z\left[\nabla_{\theta}^2g(\ystar + t\xi_{\theta}(z))\xitheta\right]\right\| \left\|\mathbb{E}_x\left[\dgx\right] + \mathbb{E}_z\left[\dgystar\right]\right\|
\end{split}
\end{equation*}
\begin{equation*}
\begin{split}
\left\langle \mathbb{E}_z\left[\nabla_{\theta}^2g(\ystar + t\xi_{\theta}(z))\xitheta\right], \mathbb{E}_x\left[\dgx\right] + \mathbb{E}_z\left[\dgystar\right]  \right\rangle \leq \\ \mathbb{E}_z\left[\left\|\nabla_{\theta}^2g(\ystar + t\xi_{\theta}(z))\xitheta\right\|\right] \left(\left\|\mathbb{E}_x\left[\dgx\right]\right\| + \left\|\mathbb{E}_z\left[\dgystar\right]\right\| \right)
\end{split}
\end{equation*}
\begin{equation*}
\begin{split}
\left\langle \mathbb{E}_z\left[\nabla_{\theta}^2g(\ystar + t\xi_{\theta}(z))\xitheta\right], \mathbb{E}_x\left[\dgx\right] + \mathbb{E}_z\left[\dgystar\right]  \right\rangle \leq \\ \mathbb{E}_z\left[\left\|\nabla_{\theta}^2g(\ystar + t\xi_{\theta}(z))\right\|\left\|\xitheta\right\|\right]\left(\sqrt{M_x} + \sqrt{M_z}\right)
\end{split}
\end{equation*}
\begin{equation}
\begin{split}
\left\langle \mathbb{E}_z\left[\nabla_{\theta}^2g(\ystar + t\xi_{\theta}(z))\xitheta\right], \mathbb{E}_x\left[\dgx\right] + \mathbb{E}_z\left[\dgystar\right]  \right\rangle \leq \\\sup_{z}\left(\left\|\nabla_{\theta}^2g(\ystar + t\xi_{\theta}(z))\right\|\right)\sup_{z}\left(\left\|\xitheta\right\|\right)\left(\sqrt{M_x} + \sqrt{M_z}\right).
\end{split}
\label{eq:upper_bound_nabla_squared}
\end{equation}
Because $\frac{\partial g(\cdot)}{\partial \theta}$ is assumed to be $L_{\theta}$-Lipschitz, it follows that $\sup_{z}\left(\left\|\nabla_{\theta}^2g(\ystar + t\xitheta)\right\| \right) \leq L_{\theta}$. Then, because for any $z$, $\|\xi_{\theta}(z)\| = \epsilon$~\eqref{eq:upper_bound_nabla_squared} becomes
\begin{equation}
\left\langle \mathbb{E}_z\left[\nabla_{\theta}^2g(\ystar + t\xi_{\theta}(z))\xitheta\right], \mathbb{E}_x\left[\dgx\right] + \mathbb{E}_z\left[\dgystar\right]  \right\rangle \leq L_{\theta}\epsilon\left(\sqrt{M_x} + \sqrt{M_z}\right).
\end{equation}
Thus, because $\epsilon \leq \frac{V\|\nabla_{\theta}L\|^2}{L_{\theta}\left(\sqrt{M_x} + \sqrt{M_z}\right)}$ for some $0 < V < 1$, 
\begin{align}
\left\langle \mathbb{E}_z\left[\nabla_{\theta}^2g(\ystar + t\xi_{\theta}(z))\xitheta\right], \mathbb{E}_x\left[\dgx\right] + \mathbb{E}_z\left[\dgystar\right]  \right\rangle &\leq L_{\theta}\left(\frac{V\|\nabla_{\theta}L\|^2}{L_{\theta}\left(\sqrt{M_x} + \sqrt{M_z}\right)}\right)\left(\sqrt{M_x} + \sqrt{M_z}\right) \\
&\leq V\|\nabla_{\theta}L\|^2.
\label{eq:upper_bound_epsilon}
\end{align}
Hence, combining ~\eqref{eq:complete_the_square} and ~\eqref{eq:upper_bound_epsilon},
\begin{align}
\left\langle \nabla_{\theta}L, \tilde{d}_{\theta} \right\rangle &\geq \|\nabla_{\theta}L\|^2 - V\|\nabla_{\theta}L\|^2 \\
&\geq (1-V)\|\nabla_{\theta}L\|^2 \\
&\geq U \|\nabla_{\theta}L\|^2.
\end{align}
\end{proof}

\subsection{Proof of Lemma ~\ref{lemma:expectation_descent}}
\begin{proof}
First, it needs to be shown that $\nabla_{\theta}L$ is Lipschitz with respect to $\theta$. Using the fact that $\frac{\partial g}{\partial \theta}$ is $L_{\theta}$-Lipschitz with respect to $\theta$, for any $\theta_{j+1},\theta_j \in \mathbb{R}^p$
\begin{align*}
\|\nabla_{\theta}L(\theta_{j+1}) - \nabla_{\theta}L(\theta_j)\| &= \left\| \mathbb{E}_x\left[\frac{\partial g(x_{\theta_{j+1}})}{\partial \theta} \right] - \mathbb{E}_z\left[\frac{\partial g(y_{\theta_{j+1}}^*(z))}{\partial \theta} \right] - \left(\mathbb{E}_x\left[\frac{\partial g(x_{\theta_{j}})}{\partial \theta} \right] - \mathbb{E}_z\left[\frac{\partial g(y_{\theta_{j}}^*(z))}{\partial \theta} \right]\right)\right\| \\
&= \left\| \mathbb{E}_x\left[\frac{\partial g(x_{\theta_{j+1}})}{\partial \theta} - \frac{\partial g(x_{\theta_{j}})}{\partial \theta}\right] + \mathbb{E}_z\left[\frac{\partial g(y_{\theta_{j}}^*(z))}{\partial \theta} -\frac{\partial g(y_{\theta_{j+1}}^*(z))}{\partial \theta}\right]\right\| \\
&\leq \left\| \mathbb{E}_x\left[\frac{\partial g(x_{\theta_{j+1}})}{\partial \theta} - \frac{\partial g(x_{\theta_{j}})}{\partial \theta}\right]\right\| + \left\|\mathbb{E}_z\left[\frac{\partial g(y_{\theta_{j}}^*(z))}{\partial \theta} -\frac{\partial g(y_{\theta_{j+1}}^*(z))}{\partial \theta}\right]\right\| \\
&\leq \mathbb{E}_x\left[ L_{\theta}\left\| \theta_{j+1} - \theta_j\right\| \right] + \mathbb{E}_z\left[L_{\theta}\left\|\theta_{j+1} - \theta_j \right\| \right] \\
&\leq L_{\theta}\left\| \theta_{j+1} - \theta_j\right\| + L_{\theta}\left\| \theta_{j+1} - \theta_j\right\| \\
&\leq 2L_{\theta} \|\theta_{j+1} - \theta_j\|.
\end{align*}
Therefore, $\nabla_{\theta}L$ is $2L_{\theta}$-Lipschitz with respect to $\theta$. \newline
Because the gradient with respect to $\theta$ of $L$, the second order Taylor series expansion of $L$ centered at $\theta=\theta_j$ satisfies, using $\theta_{j+1} = \theta_j - \alpha_j\tilde{d}_{\epsilon_j}(\theta_j)$
\begin{align*}
L(\theta_{j+1}) &\leq L(\theta_j) + \nabla_{\theta}L(\theta_j)^{\top}(\theta_{j+1} - \theta_{j}) + \frac{1}{2}(2L_{\theta})\|\theta_{j+1} - \theta_j\|^2 \\
L(\theta_{j+1}) &\leq L(\theta_j) - \alpha_j\nabla_{\theta}L(\theta_j)^{\top}\tilde{d}_{\xi_j}(\theta_j) + \alpha_j^2L_{\theta}\|\tilde{d}_{\xi_j}(\theta_j)\|^2.
\end{align*}
Taking expectation of both sides above with respect to $\xi_j$ yields
\begin{equation}
\mathbb{E}_{\xi_j}[L(\theta_{j+1})] \leq \mathbb{E}_{\xi_j}[L(\theta_j)] - \alpha_j\mathbb{E}_{\xi_j}\left[\nabla_{\theta}L(\theta_j)^{\top}\tilde{d}_{\xi_j}(\theta_j)\right] + \alpha_j^2L_{\theta}\mathbb{E}_{\xi_j}\left[\|\tilde{d}_{\xi_j}(\theta_j)\|^2\right].
\label{eq:descent_expectation_xi_j}
\end{equation}
Because $\theta_j$ depends only on $\xi_j, \xi_{j-1},...,\xi_0$, taking $\mathbb{E}_{\xi_j}[\cdot]$ only affects the LHS of \eqref{eq:descent_expectation_xi_j}. Thus, ~\eqref{eq:descent_expectation_xi_j} becomes
\begin{equation}
\mathbb{E}_{\xi_j}[L(\theta_{j+1})] \leq L(\theta_j) - \alpha_j\nabla_{\theta}L(\theta_j)^{\top}\tilde{d}_{\xi_j}(\theta_j) + \alpha_j^2L_{\theta}\|\tilde{d}_{\xi_j}(\theta_j)\|^2.
\end{equation}
Applying the result of Theorems ~\ref{theorem:lower_bound_on_inner_product} and ~\ref{theorem:norm_squared} to the above inequality,
\begin{equation}
\mathbb{E}_{\xi_j}[L(\theta_{j+1})] - L(\theta_j) \leq -\alpha_jU\|\nabla_{\theta}L(\theta_j)\|^2 + \alpha_j^2L_{\theta}\tilde{M}.
\label{eq:descent2}
\end{equation}
It can be derived that the RHS above is $\leq 0$ when $\alpha_j$ satisfies $0 < \alpha_j \leq \frac{UM_L}{L_{\theta}\tilde{M}}$:
\begin{align*}
-\alpha_jU\|\nabla_{\theta}L(\theta_j)\|^2 + \alpha_{j}^2L_{\theta}\tilde{M} &\leq 0 \\
\alpha_{j}^2L_{\theta}\tilde{M} &\leq \alpha_jU\|\nabla_{\theta}L(\theta_j)\|^2 \\
\alpha_j &\leq \frac{U\|\nabla_{\theta}L(\theta_j)\|^2}{L_{\theta}\tilde{M}} \\
\alpha &\leq \frac{UM_{L}}{L_{\theta}\tilde{M}}.
\end{align*}
\end{proof}

\textit{Proof of Theorem ~\ref{theorem:expectation_cesaro_convergence}}:
\theoremExpectationCesaroConvergence{app}
\begin{proof}
Taking the total expectation of ~\eqref{eq:descent2}, we have
\begin{equation}
\mathbb{E}[L(\theta_{j+1})] - \mathbb{E}[L(\theta_j)] \leq -\alpha_jU\mathbb{E}[\|\nabla_{\theta}L(\theta_j)\|^2] + \alpha_j^2L_{\theta}\tilde{M}.
\label{eq:conv1}
\end{equation}
Setting $j=0$ in ~\eqref{eq:conv1}, we have
\begin{align}
&\mathbb{E}[L(\theta_{1})] - \mathbb{E}[L(\theta_0)] \leq 
 -\alpha_jU\mathbb{E}[\|\nabla_{\theta}L(\theta_0)\|^2] + \alpha_j^2L_{\theta}\tilde{M} \\
&\mathbb{E}[L(\theta_{1})] \leq \mathbb{E}[L(\theta_0)]
-\alpha_jU\mathbb{E}[\|\nabla_{\theta}L(\theta_0)\|^2] + \alpha_j^2L_{\theta}\tilde{M}.
\label{eq:conv2}
\end{align}
Setting $j=1$ in ~\eqref{eq:conv1} and applying ~\eqref{eq:conv2} 
\begin{align*}
&\mathbb{E}[L(\theta_{2})] - \mathbb{E}[L(\theta_1)] \leq 
 -\alpha_jU\mathbb{E}[\|\nabla_{\theta}L(\theta_1)\|^2] + \alpha_j^2L_{\theta}\tilde{M}\\
&\mathbb{E}[L(\theta_{2})] \leq \mathbb{E}[L(\theta_1)]
-\alpha_jU\mathbb{E}[\|\nabla_{\theta}L(\theta_1)\|^2] + \alpha_j^2L_{\theta}\tilde{M}\\
&\mathbb{E}[L(\theta_{2})] \leq \mathbb{E}[L(\theta_0)]
-U\sum_{j=0}^{1}\alpha_j\mathbb{E}[\|\nabla_{\theta}L(\theta_j)\|^2] + L_{\theta}\tilde{M}\sum_{j=0}^{1}\alpha_{j}^2\\
&\mathbb{E}[L(\theta_{2})] - \mathbb{E}[L(\theta_0)]\leq 
-U\sum_{j=0}^{1}\alpha_j\mathbb{E}[\|\nabla_{\theta}L(\theta_j)\|^2] + L_{\theta}\tilde{M}\sum_{j=0}^{1}\alpha_{j}^2 \\
&\vdots  \\
&\mathbb{E}[L(\theta_{K})] - \mathbb{E}[L(\theta_0)]\leq 
-U\sum_{j=0}^{K-1}\alpha_j\mathbb{E}[\|\nabla_{\theta}J_x(\theta_j)\|^2] + L_{\theta}\tilde{M}\sum_{j=0}^{K-1}\alpha_{j}^2.
\end{align*}
Since $L(\theta)$ is bounded from below by $L_{inf}$, algebraically rearranging the last line above yields
\begin{align}
&\sum_{j=0}^{K-1}\alpha_j\mathbb{E}[\|\nabla_{\theta}L(\theta_j)\|^2] \leq 
\frac{\mathbb{E}[L(\theta_0)] - L_{inf}}{U} + \frac{L_{\theta}\tilde{M}}{U}\sum_{j=0}^{K-1}\alpha_{j}^2\\
&\frac{1}{A_K}\sum_{j=0}^{K-1}\alpha_j\mathbb{E}[\|\nabla_{\theta}L(\theta_j)\|^2] \leq \frac{\mathbb{E}[L(\theta_0)] - L_{inf}}{UA_K}+ \frac{L_{\theta}\tilde{M}}{UA_K}\sum_{j=0}^{K-1}\alpha_{j}^2.
\label{eq:conv3}
\end{align}
Using linearity of expectation, ~\eqref{eq:conv3} becomes
\begin{equation}
\mathbb{E}\left[\frac{1}{A_K}\sum_{j=0}^{K-1}\alpha_j\|\nabla_{\theta}L(\theta_j)\|^2\right] \leq \frac{\mathbb{E}[L(\theta_0)] - L_{inf}}{UA_K}+ \frac{L_{\theta}\tilde{M}}{UA_K}\sum_{j=0}^{K-1}\alpha_{j}^2.
\label{eq:conv4}
\end{equation}
Hence, since $\lim_{K \rightarrow \infty} A_K = \infty$, and $\sum_{j=0}^{\infty}\alpha_{j}^2$ converges, we have
\begin{align*}
&\lim_{K \rightarrow \infty}\mathbb{E}\left[\frac{1}{A_K}\sum_{j=0}^{K}\alpha_j\|\nabla_{\theta}J_x(\theta_j)\|^2\right] \\
&\leq \lim_{K \rightarrow \infty}\left[\frac{\mathbb{E}[L(\theta_0)]-L_{inf} + L_{\theta}\tilde{M}{\sum_{j=0}^{K-1}\alpha_{j}^2}}{U A_K}\right] \\
&= 0.
\end{align*}
\end{proof}

\subsection{Proof of Theorem ~\ref{theorem:expectation_liminf}}
\begin{proof}
Suppose, for contradiction that, for some $a > 0$
\begin{equation}
\liminf_{j \rightarrow \infty} \mathbb{E}\left[\left\| \nabla_{\theta}L(\theta_j)\right\|^2 \right] = a.
\end{equation}
Let $K \in \mathbb{N}$ and $A_K = \sum_{k=0}^{K-1}\alpha_k$. Then, using linearity of expectation
\begin{equation*}
\mathbb{E}\left[ \frac{1}{A_K}\sum_{j=0}^{K-1}\alpha_k \left\| \nabla_{\theta}L(\theta_j) \right\|^2 \right] = \frac{1}{A_K}\sum_{j=0}^{K-1}\alpha_j\mathbb{E}\left[ \left\| \nabla_{\theta}L(\theta_j) \right\|^2\right].
\end{equation*}
Taking the liminf as $K \rightarrow \infty$ on both sides of the above equation yields a contradiction, as the LHS converges to 0 but the RHS diverges. Hence, by contradiction the result follows. 
\end{proof}

\subsection{Proof of Corollary ~\ref{corollary:convergence_probability}}
\begin{proof}
This proof is similar to that of Theorem 4.11 in ~\cite{bottou2018optimization}, but we will include it here to be complete. Let $\epsilon > 0$ and let $\mathbb{E}[\cdot]$ represent total expectation. By Markov's inequality and the law of total expectation, also known as the tower property,
\begin{align}
P(\|\nabla_{\theta}L(\theta_{j(K)})\| &\geq \epsilon) = P(\|\nabla_{\theta}L(\theta_{j(K)})\|^2 \geq \epsilon^2) \\
&\leq \frac{1}{\epsilon^2}\mathbb{E}[\mathbb{E}_{j(K)}[\nabla_{\theta}L(\theta_{j(K)})]].
\label{eq:convergence_in_prob}
\end{align}
By the proof of Theorem ~\ref{theorem:expectation_cesaro_convergence} we have $\lim_{K \rightarrow \infty}\mathbb{E} \left[\sum_{j=0}^{K-1}\alpha_j \|\nabla_{\theta}L(\theta_j)\|^2 \right] < \infty$. Therefore, we must have $\lim_{j \rightarrow \infty}\mathbb{E}\left[\alpha_j\|\nabla_{\theta}L(\theta_j)\|^2\right] = 0$. Thus, by ~\eqref{eq:convergence_in_prob},
\begin{equation}
\lim_{K \rightarrow \infty} P(\|\nabla_{\theta}L(\theta_{j(K)})\| \geq \epsilon) \leq \lim_{K \rightarrow \infty } \frac{1}{\epsilon^2}\mathbb{E}[\mathbb{E}_{j(K)}[\nabla_{\theta}L(\theta_{j(K)})]] = 0.
\end{equation}
Since the choice of $\epsilon > 0$ was arbitrary, this holds $\forall \epsilon > 0$, proving convergence in probability.
\end{proof}

%%%%%%%%%%%%%%%%%%%%%%%%%%%%%%%%%%%%%%%%%%%%%%%%%%%%%%%%%%%%
% TABLE 1 — ACCURACY
%%%%%%%%%%%%%%%%%%%%%%%%%%%%%%%%%%%%%%%%%%%%%%%%%%%%%%%%%%%%

% \begin{table}[t]
% \centering
% \footnotesize
% \setlength{\tabcolsep}{3pt}  % reduce column spacing
% \renewcommand{\arraystretch}{0.95} % tighter rows
% \caption{Accuracy $\uparrow$ of the learned maps.}
% \resizebox{\linewidth}{!}{
% \begin{tabular}{lccccccccc}
% \toprule
% Datasets 
% & MUNIT & AugCG & OTDD & SinkLpL1 
% & W2 & W2,$\gamma$ 
% & FG (no $z$) & FG 
% & \textbf{Ours} \\
% \midrule
% FMNIST $\rightarrow$ MNIST 
% & 8.93 & 12.03 & 10.28 & 10.67 
% & 10.96 & 8.02 
% & 82.79 & 83.22 
% &  \\

% MNIST $\rightarrow$ MNIST-M 
% & 97.95 & 98.20 & -- & 83.26 
% & 38.77 & 37.00 
% & 95.27 & 94.62 
% &  \\
% \bottomrule
% \end{tabular}
% }
% \end{table}


%%%%%%%%%%%%%%%%%%%%%%%%%%%%%%%%%%%%%%%%%%%%%%%%%%%%%%%%%%%%
% TABLE 2 — FID
%%%%%%%%%%%%%%%%%%%%%%%%%%%%%%%%%%%%%%%%%%%%%%%%%%%%%%%%%%%%

\section{Implementation Details}\label{appen:implementation_details}

\subsection{Datasets}

\subsubsection{Synthetic Gaussians}
For the synthetic high-dimensional Gaussian experiments in Section~\ref{sec:high_dim_gaussians_rewrite}, we use the data generation procedure of \cite{park2025neural}, which follows the standard protocol of \cite{korotin2021neural} for constructing Gaussian pairs with controlled covariance spectra.

In this setup, each covariance matrix is defined as
\[
\Sigma = Q \Lambda Q^\top,
\]
where $Q$ is an orthogonal matrix obtained via QR decomposition of a Gaussian random matrix, and $\Lambda$ contains eigenvalues sampled from a log-uniform distribution,
\[
\lambda_i \sim \exp(\mathcal{U}(a,b)).
\]
This yields well-conditioned but heterogeneous Gaussian covariances.

Given two Gaussian distributions $\mathcal{N}(0, \Sigma_0)$ and $\mathcal{N}(0, \Sigma_1)$, the optimal transport map is linear, $T(x)=\Gamma x$, where
\[
\Gamma = \Sigma_0^{-1/2}
\left(\Sigma_0^{1/2} \Sigma_1 \Sigma_0^{1/2}\right)^{1/2}
\Sigma_0^{-1/2}.
\]

Samples from $\mathcal{N}(0, \Sigma_0)$ and $\mathcal{N}(0, \Sigma_1)$ are drawn following the same procedure as in \cite{park2025neural}, with $n=10^5$ samples for both training and evaluation. Independent test samples are used to evaluate transport accuracy, where the ground-truth map $T(x)=\Gamma x$ is applied to source samples, and $\Gamma^{-1}$ is used for inverse-consistency checks.

\subsubsection{UCI Physics Data}
The UCI physics dataset is used following the preprocessing protocol of \cite{onken2021ot}, and we adopt the same train/validation/test split and normalization procedure as in the original work.

\paragraph{GAS Dataset}
For the GAS dataset, %we follow the preprocessing pipeline of \cite{bengio2011gas}. 
the raw data is loaded from \texttt{ethylene\_CO.csv}, and the variables \texttt{Meth}, \texttt{Eth}, and \texttt{Time} are removed. Highly correlated features are iteratively removed by discarding variables whose pairwise correlation exceeds $0.98$. After this cleaning step, each feature is standardized using z-score normalization based on the training data statistics. The dataset is then split into training, validation, and test sets using a $80/10/10$ split.

\paragraph{POWER Dataset}
For the POWER dataset, we use the preprocessed data from \texttt{data.npy}. Two features are removed as in the original preprocessing pipeline, and additional synthetic noise is added to several variables following the procedure of prior work. The dataset is then randomly shuffled and split into training, validation, and test sets using a $80/10/10$ ratio.

\paragraph{HEPMASS Dataset}
For the HEPMASS dataset, we follow the preprocessing pipeline of prior work on tabular density estimation. We use the split provided in \texttt{1000\_train.csv} and \texttt{1000\_test.csv}, and restrict the data to the positive class only by removing background noise samples (class label $0$). In addition, the label column and one redundant feature column in the test set are removed.

We consider three normalization schemes: static normalization, min-max scaling, and SVD-based whitening. In the experiments reported in the main paper, we use the \textit{min-max scaling} variant (denoted as \texttt{scale}), where each feature is normalized as
\[
x' = \frac{x - \mu}{s}, \quad
\mu = \frac{\max(x) + \min(x)}{2}, \quad
s = \frac{\max(x) - \min(x)}{2}.
\]
Statistics are computed using the training set only. The resulting transformation maps each feature approximately to $[-1, 1]$. We further remove features with highly repetitive values following the procedure in the original preprocessing code.

\paragraph{MINIBOONE Dataset}
For the MINIBOONE dataset, we use the preprocessed version provided as a NumPy array. The dataset is split into training, validation, and test sets using an $80/10/10$ ratio.

We apply min-max scaling (denoted as \texttt{scale}) using statistics computed on the combined training and validation set:
\[
x' = \frac{x - \mu}{s}, \quad
\mu = \frac{\max(x) + \min(x)}{2}, \quad
s = \frac{\max(x) - \min(x)}{2}.
\]
This normalization is applied independently to each feature, resulting in an approximate $[-1, 1]$ range per dimension.

\paragraph{BSDS300 Dataset}
The BSDS300 dataset consists of natural image patches extracted from the Berkeley Segmentation Dataset. We use the standard train/validation/test split provided in the preprocessed HDF5 file \texttt{BSDS300.hdf5}, following prior work on image-based optimal transport and density modeling.

Each sample corresponds to a vectorized image patch, and no additional feature engineering is applied. We directly load the precomputed splits for training, validation, and testing.

Since the data consists of image patches, we additionally record the corresponding spatial resolution, given by
\[
\text{image size} = \left[\sqrt{d+1}, \sqrt{d+1}\right],
\]
where $d$ denotes the dimensionality of the vectorized patch. No further normalization beyond the provided preprocessing is applied.

\subsubsection{Synthetic Class-Conditional OT Datasets}
We construct three synthetic benchmarks for class-conditional optimal transport using mixtures of two-dimensional Gaussian clusters arranged on a circular manifold. All datasets are generated on a ring of radius $r=0.9$ with isotropic Gaussian noise $\epsilon = 0.07$. Each class corresponds to a subset of Gaussian components placed at fixed angular positions.

\paragraph{Crossed Ring Gaussians}
In the Crossed Ring Gaussian dataset, each distribution consists of two Gaussian modes placed at distinct angular positions on the unit circle. Specifically, the source distribution is defined by two clusters centered at angles $0$ and $\pi/2$, while the target distribution is defined by clusters at angles $\pi/4$ and $3\pi/4$. This creates a crossed transport structure between classes. The total number of samples is $N = 50{,}000$, with equal allocation per mode.

\paragraph{Horizontal Swapped Gaussians}
We construct an additional synthetic benchmark consisting of two one-dimensional Gaussian clusters embedded in $\mathbb{R}^2$. The dataset is designed to evaluate class-conditional optimal transport under near-degenerate geometric structure.

The source distribution consists of two Gaussian components centered at $(-0.23, 0)$ and $(0.23, 0)$, while the target distribution consists of two components centered at $(0.7, 0)$ and $(-0.7, 0)$, respectively. This induces a horizontal swapping structure along the $x$-axis.

Formally, each component is generated as
\[
x = \mu + \epsilon \cdot \mathcal{N}(0, I_2),
\]
where $\mu \in \mathbb{R}^2$ denotes the cluster center and $\epsilon = 0.05$ controls the noise level. The total number of samples is $N = 50{,}000$, equally split across components.

This dataset tests the ability of class-conditional OT methods to recover correct pairwise matching in a low-dimensional, highly structured setting.
\paragraph{Four-Mode Ring Gaussians}
The Four-Mode Ring Gaussian dataset consists of four Gaussian components per distribution placed uniformly on the unit circle. The source distribution uses angles $\{0, \pi/2, \pi, 3\pi/2\}$, while the target distribution uses $\{\pi/4, 3\pi/4, 5\pi/4, 7\pi/4\}$. This results in a dense multimodal matching problem with $N = 5{,}000$ total samples.


For all datasets, each mode is sampled as
\[
x = r \cdot (\cos\theta, \sin\theta) + \epsilon \cdot \mathcal{N}(0, I_2),
\]
where $r=0.9$ and $\epsilon=0.07$. Samples are concatenated across modes to form the full empirical distributions. The resulting datasets are used to evaluate class-conditional optimal transport under structured multimodal alignments.

\subsubsection{Image Data}
Normalization is performed using a min-max scaling transformation computed from the combined training and validation set:
\[
x' = \frac{x - \mu}{s}, \quad
\mu = \frac{\max(x) + \min(x)}{2}, \quad
s = \frac{\max(x) - \min(x)}{2}.
\]
This maps each feature approximately into the range $[-1, 1]$.
All images are resized to $32 \times 32$ pixels and normalized to $[0, 1]$.
For training the OT map, we construct a \textbf{class-paired dataset} by randomly
drawing matched $(x_{\mathrm{src}}, x_{\mathrm{tgt}}, k)$ triples where both images
belong to the same class $k \in \{0, \ldots, 9\}$, drawn uniformly at random from
FashionMNIST (source) and MNIST (target) training sets (60{,}000 images each).

\subsection{Kantorovich Potential Network}
For the high-dimensional Gaussian experiments in Section~\ref{sec:high_dim_gaussians_rewrite}, we employ the DenseICNN architecture, a fully connected neural network augmented with input-quadratic skip connections, in order to ensure a fair comparison with the baseline models of \cite{korotin2021neural}. When implementing our method and the NCF baseline, we remove the nonlinearity constraints that are typically imposed to enforce convexity.
Following \cite{korotin2021neural}, we adopt the network architecture DenseICNN[1; max(2$d$,64), max(2$d$,64), max($d$,32)] for a $d$-dimensional problem. The model is optimized using the Adam optimizer with a fixed learning rate of $10^{-4}$, independent of the input dimension.

The same network architecture is also used for the 2D class-conditional optimal transport experiments in Section~\ref{sec:uci_physics}, Section~\ref{sec:ccot_2d}, and ablation studies in Section \ref{sec:ablation}.

For the image translation experiments in Section \ref{sec:mnist_fmnist}, the potential $g_\theta(x, k)$ is parameterized by a scalar-valued network
taking as input the concatenation $[x;\, \mathbf{c}_k] \in \mathbb{R}^{25}$,
where $\mathbf{c}_k \in \{0,1\}^{10}$ is the class one-hot vector.
The network decomposes as
\begin{equation}
  g_\theta(x, k)
  = \underbrace{\tfrac{1}{2}\, x^\top (A^\top\! A)\, x + b^\top x + c}_{%
      \text{quadratic}}
  + \underbrace{w^\top \Phi(x, k)}_{\text{nonlinear}},
\end{equation}
where $A \in \mathbb{R}^{25 \times 25}$ is Xavier-initialized, enforcing a
positive semi-definite quadratic component.
The nonlinear branch $\Phi$ is a residual network: a linear opening layer
$\mathbb{R}^{25} \to \mathbb{R}^{512}$, followed by four hidden layers of
width $512$ with residual connections weighted by $h = \tfrac{1}{4}$, and a
linear readout. The activation function is softplus throughout.

\subsection{Implicit Fixed-Point Solver}

The proximal fixed point
$y^\star(z, k) = \operatorname{prox}_{g_\theta}(z, k)$ is computed by iterating
\begin{equation}
  y^{(t+1)}
  = y^{(t)} - \alpha\!\left(y^{(t)} - z + \nabla_y g_\theta(y^{(t)}, k)\right),
\end{equation}
with step size $\alpha = \min\!\bigl(0.1/\hat{L},\; 10^{-2}\bigr)$, where
$\hat{L}$ is the spectral norm of the Hessian of $g_\theta$ estimated via
20 steps of power iteration.
Iterations run until
$\|\nabla_y g_\theta(y^\star, k) + y^\star - z\|_\infty < 10^{-3}$
or for at most $10^4$ steps.
The previous batch's fixed point is used as a warm start.

\subsection{OT Map Training}

The OT map is trained to minimize
\begin{equation}
  \mathcal{L}(\theta)
  = \mathbb{E}_{\mu_k}\!\left[g_\theta(x,k)\right]
  - \mathbb{E}_{\nu_k}\!\left[g_\theta\!\left(y^\star(z,k),\, k\right)\right]
  + \lambda_{\mathrm{MMD}}\,\widehat{\mathrm{MMD}}^2\!\left(T(x),\, z\right),
\end{equation}
where $T(x) = x + \nabla_x g_\theta(x, k)$ is the forward transport map
and $\widehat{\mathrm{MMD}}^2$ is the unbiased multi-scale RBF kernel
estimator. We use five kernel bandwidths,
$\sigma^2 \in \{0.25,\, 0.5,\, 1.0,\, 2.0,\, 4.0\} \times \sigma^2_{\mathrm{median}}$, which are only applied in the image-based experiments in Section~\ref{sec:mnist_fmnist}. Consequently, we set $\lambda_{\mathrm{MMD}} = 0$ for the experiments in Sections~\ref{sec:high_dim_gaussians_rewrite}, \ref{sec:uci_physics}, \ref{sec:ccot_2d}, and \ref{sec:ablation}.

\begin{table}[h]
\centering
\begin{tabular}{lc}
\toprule
\textbf{Hyperparameter} & \textbf{Value} \\
\midrule
Epochs                        & 20 \\
Batch size                    & 64 \\
MMD subsample size            & 5000 \\
$\lambda_{\mathrm{MMD}}$      & $10^{-3}$ \\
Fixed-point tolerance $\tau$  & $10^{-3}$ \\
Max fixed-point iterations    & $10^4$ \\
Optimizer                     & Adam \\
Learning rate                 & $10^{-4}$ \\
LR schedule                   & ReduceLROnPlateau (factor $0.5$, patience $5$) \\
Minimum LR                    & $10^{-7}$ \\
Floating-point precision      & float64 \\
\bottomrule
\end{tabular}
\caption{OT map training hyperparameters.}
\end{table}




\subsection{Conditional VAE for Image Task}

\paragraph{Architecture}

Both the FashionMNIST and MNIST VAEs share identical architecture.
The \textbf{encoder} consists of four strided convolutional layers
(kernel $4\times4$, stride $2$, padding $1$), with channel widths
$1 \to 64 \to 128 \to 256 \to 512$, each followed by BatchNorm and
LeakyReLU($0.2$).
The flattened $512\times2\times2$ feature map is concatenated with the
class one-hot vector and projected to the mean $\mu$ and log-variance
$\log\sigma^2$ of the latent distribution via linear layers, yielding a
latent dimension of $d = 15$.

The {decoder} takes the concatenation of the latent code and class
one-hot as input, projects it to a $512\times2\times2$ feature map, and
upsamples through four transposed convolutional layers (kernel $4\times4$,
stride $2$) with channel widths $512 \to 256 \to 128 \to 64 \to 1$.
Each upsampling block is followed by GroupNorm(16) and three residual blocks
(two $3\times3$ convolutions with GroupNorm and a skip connection).
The final output is passed through a Sigmoid activation.
GroupNorm is used in the decoder (instead of BatchNorm) to avoid
train/eval distribution mismatch.

% The VAE is trained with a $\beta$-VAE objective:
% \begin{equation}
%   \mathcal{L}_{\mathrm{VAE}}
%   = \mathbb{E}\bigl[\log p(x \mid z)\bigr]
%   - \beta\, D_{\mathrm{KL}}\!\bigl(q(z\mid x) \;\|\; \mathcal{N}(0,I)\bigr),
% \end{equation}
% where the reconstruction term uses binary cross-entropy.

\paragraph{VAE Training Hyperparameters}

\begin{table}[h]
\centering
\begin{tabular}{lcc}
\toprule
\textbf{Hyperparameter} & \textbf{FashionMNIST VAE} & \textbf{MNIST VAE} \\
\midrule
Latent dimension $d$      & 15   & 15  \\
$\beta_{\max}$            & 1.0  & 0.1 \\
KL warmup epochs          & 20 (linear ramp) & 0 \\
Epochs                    & 1000 & 1000 \\
Batch size                & 128  & 128  \\
Optimizer                 & Adam & Adam \\
Learning rate             & $10^{-3}$ & $10^{-3}$ \\
LR schedule               & \multicolumn{2}{c}{ReduceLROnPlateau (factor $0.5$, patience $100$)} \\
Minimum LR                & \multicolumn{2}{c}{$10^{-7}$} \\
\bottomrule
\end{tabular}
\caption{VAE training hyperparameters.}
\end{table}
Both VAEs are frozen after pre-training; their parameters are not updated
during OT map training.



% \section*{Acknowledgments}


\bibliographystyle{siamplain}
\bibliography{mybib}

% \input{ex_supplement}

\end{document}
